\documentclass{article}

\usepackage{amsmath}
\usepackage{amssymb}
\usepackage{amsthm}
\usepackage[margin=1in]{geometry}
\usepackage[hidelinks]{hyperref}
\newcommand{\proofheading}[1]{%
    \par\addvspace{\smallskipamount}%
    \noindent\textbf{#1}\par\nobreak
}

\DeclareMathOperator{\ord}{ord}
\newtheorem{proposition}{Proposition}[section]
\newtheorem{theorem}[proposition]{Theorem}
\newtheorem{lemma}[proposition]{Lemma}
\newtheorem{question}[proposition]{Question}
\theoremstyle{remark}
\newtheorem*{remark}{Remark}

\title{A Hereditarily Atomic Domain without ACCP}
\author{Matthew Toeniskoetter}

\begin{document}

\maketitle

\begin{abstract}
    An integral domain for which every subring is atomic is said to be hereditarily atomic.
    Every hereditarily atomic cancellative commutative monoid satisfies the ascending chain condition on principal ideals (ACCP), but it remained an open question whether every hereditarily atomic domain satisfies ACCP.
    We construct a hereditarily atomic domain which does not satisfy ACCP, answering this question in the negative.
    This domain is an overring of the ring of integer-valued polynomials $\operatorname{Int} (\mathbb{Z})$, and its subrings that fail ACCP are precisely those of finite additive index.
\end{abstract}

\section{Introduction}

A nonzero nonunit $a$ in an integral domain is an \emph{atom} if $a = bc$ forces $b$ or $c$ to be a unit, and a domain is \emph{atomic} if every nonzero nonunit is a finite product of atoms.
A domain satisfies the \emph{ascending chain condition on principal ideals} (ACCP) if every ascending chain of principal ideals eventually stabilizes.
It is easily shown that ACCP implies atomicity, but in 1974 Grams showed that the converse fails by constructing an atomic domain with an infinite strictly ascending chain of principal ideals \cite[Theorem~1.3]{Grams}.

Throughout, subrings contain the identity of the ambient ring.
A domain is \emph{hereditarily atomic} if every one of its subrings is atomic.
Coykendall, Gotti, and Hasenauer conjectured that a domain is hereditarily atomic if and only if every one of its subrings satisfies ACCP \cite[Conjecture~6.1]{CGH}.
Gotti and Li proved the analogous conjecture for commutative cancellative monoids by showing that a monoid failing ACCP contains a submonoid in which some nonunit has no factorization into atoms \cite[Theorem~4.2 and Corollary~4.3]{GLmonoids}.
However, passing from a multiplicative submonoid of a domain to the subring it generates can introduce new elements through sums and differences, which can then yield new factorizations.
Theorem~\ref{thm:counterexample} refutes the domain conjecture in characteristic zero.

\begin{theorem}
    \label{thm:counterexample}
    There exists a hereditarily atomic integral domain $D$ with fraction field $\mathbb{Q}(t)$ that fails ACCP.
\end{theorem}

Our domain $D$ is an overring of the ring $\operatorname{Int}(\mathbb{Z})$ of integer-valued polynomials on $\mathbb{Z}$.
Following the valuation description of Chabert and Peruginelli, we impose integrality at two $p$-adic points over each prime and allow poles at the roots of selected irreducible polynomials \cite[\S4]{CPOverrings}.
Chang constructs almost Dedekind domains between $\mathbb{Z}[t]$ and $\mathbb{Q}[t]$ with prescribed ideal class groups \cite[Theorem~3.5]{Chang}, and Peruginelli uses $p$-adic evaluation to realize such groups in Dedekind domains \cite[Theorem~3.1]{Peruginelli}.
We choose the irreducible factors allowed in our denominators by combining coefficient congruences with Eisenstein's criterion at an auxiliary prime, as Frisch does to obtain irreducible polynomials while preserving the fixed divisors of their products \cite[Lemma~6]{Frisch}.
We prescribe the valuations so that we can divide $t^2$ successively by nonunits of $D$.
The resulting quotients generate an infinite strictly ascending chain of principal ideals beginning at $(t^2)$.
To obtain atomic factorizations in every subring $T \subseteq D$, we impose further interpolation conditions on the same polynomials.
A \textbf{decomposition} of a nonzero $x \in T$ is an expression $x = y_1 \cdots y_\ell$ with each $y_j$ a nonunit of $T$.
It is a factorization into atoms of $T$ when every $y_j$ is an atom of $T$.
If a nonzero element of $T$ admits arbitrarily long decompositions into nonunits of $T$, these conditions force $mD \subseteq T$ for some positive integer $m$.
We construct atoms $a_n$ with $a_n \equiv 1 \pmod{mD}$ for all sufficiently large $n$, so whenever such an $a_n$ divides $x \in T$ in $D$, both $a_n$ and $x/a_n$ belong to $T$.
Dividing by a sufficiently late $a_n$ lowers the order at $t = 0$ by one, until we reach an element of order zero or one, where decomposition lengths are bounded.

The same conductor argument lets us determine exactly which subrings of $D$ fail ACCP.
Recall that an atomic domain is a \emph{bounded factorization domain} (BFD) if, for each nonzero nonunit, the lengths of its factorizations into atoms are bounded.
Write $[D : T]$ for the index of the additive group of $T$ in that of $D$.
The \emph{conductor} $(T : D)$ is the largest ideal of $D$ contained in $T$.

\begin{theorem}
    \label{thm:subring-classification}
    For a subring $T$ of the domain $D$ constructed in Theorem~\ref{thm:counterexample}, the following are equivalent.
    \begin{enumerate}
        \item $T$ is not ACCP.
        \item $T$ is not a BFD.
        \item $[D : T]$ is finite.
        \item $T = \mathbb{Z} + nD$ for some integer $n \ge 1$.
    \end{enumerate}
    In this case $n = [D : T]$ and $(T : D) = nD$.
    If $[D : T]$ is infinite, every subring of $T$ is a bounded factorization domain.
\end{theorem}

The original AI-authored account of this construction and the Lean formalization of Theorems~\ref{thm:counterexample} and~\ref{thm:subring-classification} are available together in \cite{LeanHA}.

\section{Background on valuations and divisibility}
\label{sec:background}

For $f = P/Q \in \mathbb{Q}(t)^\times$ in lowest terms and a monic irreducible $h \in \mathbb{Q}[t]$, write $\ord_h(f)$ for the exponent of $h$ in $P$ minus its exponent in $Q$.
In particular, $\ord_t(f)$ is the order of vanishing at $t = 0$.
The finite poles of $f$ are the roots of $Q$ in $\overline{\mathbb{Q}}$, and $f$ has a pole at infinity precisely when $\deg P > \deg Q$.

Let $\mathbb{Q}_p$ denote the $p$-adic numbers and $\mathbb{Z}_p$ denote the $p$-adic integers.
Normalize the $p$-adic valuation on $\mathbb{Q}_p$ by $v_p(p) = 1$ and write $v_p(0) = +\infty$.
For $z \in \mathbb{Q}_p$, we have $z \in \mathbb{Z}_p$ if and only if $v_p(z) \ge 0$, and $z \in \mathbb{Z}_p^\times$ if and only if $v_p(z) = 0$.
If $\gamma \in \mathbb{Z}_p$ is transcendental over $\mathbb{Q}$, evaluation at $\gamma$ embeds $\mathbb{Q}(t)$ in $\mathbb{Q}_p$ and gives the valuation $f \mapsto v_p(f(\gamma))$.
Its value on $t$ is $v_p(\gamma)$, whereas $\ord_t(t) = 1$.

Let $R$ be the intersection of the valuation rings of a family $\mathcal{V}$ of valuations on a field $K$.
For nonzero $x, y \in R$, testing $y/x \in R$ gives
\[
    x \mid y \text{ in } R
    \quad\Longleftrightarrow\quad
    v(x) \le v(y) \text{ for every } v \in \mathcal{V}.
\]

\section{The domain and the prescribed valuations}

We work in the function field $K = \mathbb{Q}(t)$.
To construct $D$, we first specify the polynomial factors which may occur in the denominators of rational functions.
Let $F_1, F_2, \ldots$ be distinct monic irreducible polynomials in $\mathbb{Q}[t]$, none of which are $t$ itself, and define
\[
    B = \mathbb{Q}[t, F_1^{-1}, F_2^{-1}, \ldots].
\]
Thus $B$ consists of the rational functions that can be written with denominator a finite product of powers of the $F_i$.
Every element of $B$ is regular at $t = 0$, and $1/t$ does not belong to $B$.

Next, we impose two integrality conditions at each rational prime $p$.
For each $p$, we choose two points $\alpha_p, \beta_p \in \mathbb{Z}_p$, both transcendental over $\mathbb{Q}$.
The \textbf{primary point} $\alpha_p$ may be any such point satisfying $v_p(\alpha_p) = 1$ for odd $p$ and $v_2(\alpha_2) = 2$ for $p = 2$.
The \textbf{guard point} $\beta_p$ is a unit, so $v_p(\beta_p) = 0$.
Section~\ref{sec:realization} constructs the polynomials $F_i$ and specifies how to choose the residue of each $\beta_p$ modulo $p$.
Once that residue is fixed, any transcendental lift will do.

We now define
\[
    D = \{ f \in B \mid f(\alpha_p), f(\beta_p) \in \mathbb{Z}_p
                        \text{ for every rational prime } p \}.
\]
Since these integrality conditions are preserved under addition, subtraction, and multiplication, it follows that $D$ is a subring of $B$.

Write $E_p = \{\alpha_p, \beta_p\}$ and let $A$ denote the polynomial part of $D$,
\[
    A = D \cap \mathbb{Q}[t]
      = \{ f \in \mathbb{Q}[t] \mid f(E_p) \subseteq \mathbb{Z}_p \text{ for every prime } p \}.
\]
Every $f \in \operatorname{Int}(\mathbb{Z})$ maps $\mathbb{Z}_p$ into $\mathbb{Z}_p$ by continuity, since $\mathbb{Z}$ is dense in $\mathbb{Z}_p$.
Thus $\mathbb{Z}[t] \subseteq \operatorname{Int}(\mathbb{Z}) \subseteq A \subseteq D$.
In the classification of Chabert and Peruginelli, $A$ is the polynomial overring corresponding to the closed sets $E_p$, written $\operatorname{Int}_{\mathbb{Q}}(\prod_p E_p, \widehat{\mathbb{Z}})$ with $\widehat{\mathbb{Z}} = \prod_p \mathbb{Z}_p$ \cite[Theorem~6.2]{CPOverrings}.
For example, if $r \in \mathbb{Z}$ represents $\beta_3$ modulo $3$, then $t(t-r)/3$ belongs to $A$, but its value at an integer in the third residue class modulo $3$ is not an integer.

We describe a membership test for $B$ and for $D$ in terms of valuations.
A nonzero rational function $f \in K$ belongs to $B$ if and only if every monic irreducible factor of its reduced denominator is one of the $F_i$.
Equivalently, $f \in B$ if and only if $\ord_h(f) \ge 0$ for every monic irreducible $h \in \mathbb{Q}[t]$ other than the $F_i$.
Thus $D$ is the intersection of the rings $\mathbb{Q}[t]_{(h)}$ for these $h$ and the evaluation valuation rings at the points $\alpha_p$ and $\beta_p$, as in \cite[\S4]{CPOverrings}.
An element $f \in B$ belongs to $D$ precisely when $v_p(f(\alpha_p)) \ge 0$ and $v_p(f(\beta_p)) \ge 0$ at every prime $p$.
For $f \in K^\times$, its \textbf{valuation profile} has a polynomial coordinate $\ord_h(f)$ for each monic irreducible $h \in \mathbb{Q}[t]$ other than the $F_i$, and a primary coordinate $v_p(f(\alpha_p))$ and a guard coordinate $v_p(f(\beta_p))$ for each rational prime $p$.
By Section~\ref{sec:background}, divisibility in $D$ is determined by coordinatewise comparison of these profiles.

Consider the profile of $t$.
We have $\operatorname{ord}_t(t) = 1$, and $\operatorname{ord}_h(t) = 0$ for every other monic irreducible $h$.
At each guard the value is $v_p(\beta_p) = 0$.
Write $\mathbf{v}(f)$ for the \textbf{primary profile} of $f$, the sequence of its primary values $v_p(f(\alpha_p))$ in increasing order of the primes $2, 3, 5, \ldots$.
Its support is the set of indices at which its value is nonzero.
Evaluating $t$ at the primary point gives $\alpha_p$, and squaring doubles each valuation, so
\[
    \mathbf{v}(t) = (2, 1, 1, 1, \ldots),
    \qquad
    \mathbf{v}(t^2) = (4, 2, 2, 2, \ldots).
\]

We can now prescribe a divisor of $t^2$.
Section~\ref{sec:realization} will construct an element $q_0$ with
\[
    \mathbf{v}(q_0) = (2, 1, 2, 0, 0, \ldots)
\]
and with every guard and polynomial coordinate equal to zero.
Comparing profiles gives $q_0 \mid t^2$.
At the primary over $5$, however, $v_5(q_0(\alpha_5)) = 2 > 1 = v_5(\alpha_5)$, so $q_0 \nmid t$.

Index the primes in increasing order as $p_0 = 2$, $p_1 = 3$, $p_2 = 5$, and so on.
For each $n \ge 0$, Section~\ref{sec:realization} will construct $q_n$ with primary values $2, 1, 2$ at $p_{2n}, p_{2n + 1}, p_{2n + 2}$, respectively, and value zero at every other coordinate of its profile.
Since the nonunits $q_0, q_2, q_4, \ldots$ have pairwise disjoint primary supports, dividing $t^2$ successively by them gives $(t^2) \subsetneq (t^2/q_0) \subsetneq (t^2/(q_0 q_2)) \subsetneq \cdots$, so $D$ fails ACCP.

To construct the atoms we will use later, consider what happens when we divide $t$ by $q_0$.
Subtracting their primary values gives
\[
    \mathbf{v}(t/q_0) = (0, 0, -1, 1, 1, \ldots).
\]
The only negative entry is at the primary over $5$.
We want the primary profile of $t$ shifted two places to the right, so the entry at $5$ must increase from $-1$ to $2$.
We therefore prescribe $b_1$ with $v_5(b_1(\alpha_5)) = 3$ and value zero at every other coordinate.
Let $a_1 = t b_1/q_0$.
Multiplication by $b_1$ gives
\[
    \mathbf{v}(a_1) = (0, 0, 2, 1, 1, \ldots).
\]
The guard and polynomial coordinates of $a_1$ agree with those of $t$, since these coordinates are zero for both $b_1$ and $q_0$.
Thus $a_1 \in D$ and $\ord_t(a_1) = 1$.

For each $n \ge 0$, we prescribe the following primary values for $b_n$ and $q_n$, with every other coordinate of both profiles equal to zero.
\begin{equation*}
    \tag{P}\label{eq:prescribed-profiles}
    \begin{array}{c|ccc}
        & p_{2n} & p_{2n + 1} & p_{2n + 2} \\
        \hline
        b_n & 3 & 0 & 0 \\
        q_n & 2 & 1 & 2
    \end{array}
\end{equation*}
In particular, $b_0$ has $v_2(b_0(\alpha_2)) = 3$.
Let $a_0 = t$ and define
\[
    a_{n + 1} = \frac{a_n b_{n + 1}}{q_n}.
\]
On the three primary coordinates beginning at $p_{2n}$, this step changes $(2, 1, 1)$ to $(2, 1, 1) + (0, 0, 3) - (2, 1, 2) = (0, 0, 2)$.
Thus each step shifts the initial $2$ two places to the right, leaving zeros before it and ones after it.
The guard and polynomial values agree with those of $t$ throughout, so each $a_n$ belongs to $D$ and satisfies $\ord_t(a_n) = 1$.
We will prove that the $a_n$ are atoms of $D$ and use them to factor elements in subrings that need not contain $t$.

Section~\ref{sec:realization} will construct the two families in the order
\[
    b_0,\ b_1,\ q_0,\ b_2,\ q_1,\ b_3,\ q_2,\ \ldots.
\]
When we choose $q_n$, both $a_n$ and $b_{n + 1}$ are already known, so the numerator in $a_{n + 1} = a_n b_{n + 1}/q_n$ is fixed.
At each prescribed evaluation where the numerator and denominator must be units, we can require $q_n$ to agree with the numerator modulo a power of $p$, making $a_{n + 1}$ congruent to $1$ modulo that power at the same point.
We will use this freedom to ensure that for each fixed integer $m \ge 2$, we have $a_n \equiv 1 \pmod{mD}$ for all sufficiently large $n$.

\section{Realizing the prescribed valuations}
\label{sec:realization}

Let $U_i$ be the $i$th term of the sequence $b_0, b_1, q_0, b_2, q_1, \ldots$, starting with $i = 1$.
At the $i$th stage, the earlier polynomials $F_1, \ldots, F_{i - 1}$ and the evaluation points already assigned are fixed.
We choose the denominator $B_i$ to be $(F_1 \cdots F_{i - 1})^{K_i}$ for some $K_i > 0$, and construct a new monic irreducible polynomial $F_i \in \mathbb{Q}[t]$ with $F_i(0) \ne 0$.
Then the rational function constructed at this stage is $U_i = F_i/B_i$.
In particular, $B_1 = 1$ and $U_1 = F_1 = b_0$.

The profiles of the individual $U_i$ explain the failure of ACCP.
To prove atomicity in an arbitrary subring, we also need to control products of the $U_i$ that belong to that subring.
At the $i$th stage, consider the finite list of products of integer powers of $U_1, \ldots, U_i$,
\[
    \mathcal{W}_i =
    \left\{
        U_i^a \prod_{h = 1}^{i - 1} U_h^{k_h}
        \;\middle|\;
        1 \le a \le i,\ -i \le k_h \le i,\ a, k_h \in \mathbb{Z}
    \right\}.
\]
In $\mathcal{W}_i$, the newest factor has positive exponent, while the earlier factors occur in either the numerator or the denominator.
By choosing the power in $B_i$ large enough in $U_i = F_i / B_i$, we can keep every earlier $F_h$ in the denominator of every product on this finite list.

We say a point $\lambda \in \overline{\mathbb{Q}}$ is a \textbf{retained point} if $F_j(\lambda) \ne 0$ for every $j \ge 1$.
Since $U_j$ is a ratio of $F_j$ with a product of powers of $F_1, \ldots, F_{j-1}$, each $U_j$ can be evaluated at retained points.
Thus every product in each list $\mathcal{W}_i$ across all $i \ge 1$ can be evaluated at retained points.

We write $w'$ for the formal derivative of $w$ with respect to $t$.

\begin{proposition}[Simultaneous realization]
    \label{prop:simultaneous-realization}
    The polynomials $F_i$, the denominators $B_i$, and the evaluation points $\alpha_p, \beta_p$ can be chosen to satisfy all of the following conditions.
    \begin{enumerate}
        \item The valuation profiles of $b_n$ and $q_n$ are those prescribed in \eqref{eq:prescribed-profiles}.
        Moreover, for each nonzero $f \in K$ and all sufficiently large primes $p$,
        \[
            v_p(f(\alpha_p)) = \ord_t(f), \qquad v_p(f(\beta_p)) = 0.
        \]
        The recursively defined elements $a_n$ belong to $D$, and for each integer $m \ge 2$ we have $a_n \equiv 1 \pmod{mD}$ for all sufficiently large $n$.

        \item Each $w \in \mathcal{W}_i$ with exponent $a$ on $U_i$ has the form $w = F_i^a/H_{i, w}$, where the denominator $H_{i, w}$ is a monic product of positive powers of all the earlier polynomials $F_1, \ldots, F_{i - 1}$, determined by the exponent tuple and $B_i$ before $F_i$ is chosen.
        For $i = 1$ this is the empty product, so $H_{1, w} = 1$.
        We choose $\deg F_i$ large enough that $a \deg F_i > \deg H_{i, w}$ for every product on the list.
        Therefore, the poles of each $w \in \mathcal{W}_i$ on $\mathbb{P}^1(\overline{\mathbb{Q}})$ are precisely $\infty$ and the roots of the earlier polynomials $F_h$ with $h < i$.

        \item The lists $\mathcal{W}_i$ satisfy the following separation and derivative conditions.
        \begin{enumerate}
            \item For each pair of distinct retained points $\lambda, \mu$, all sufficiently large $i$ satisfy $w(\lambda) \ne w(\mu)$ for every $w \in \mathcal{W}_i$.
            \item For each retained point $\lambda$, all sufficiently large $i$ satisfy $w'(\lambda) \ne 0$ for every $w \in \mathcal{W}_i$.
            \item For each prime $p$ and each $\gamma \in \{\alpha_p, \beta_p\}$,
            \[
                \max_{w \in \mathcal{W}_i} v_p(w'(\gamma)) \longrightarrow -\infty
                \qquad \text{as } i \longrightarrow \infty.
            \]
        \end{enumerate}

        \item Fix a prime $p$ and an integer $s \ge 0$.
        Define
        \[
            A_s = \mathbb{Z}_p[p^s t] \subseteq \mathbb{Q}_p(t), \qquad g = (p^s t - p^s\alpha_p)(p^s t - p^s\beta_p) \in A_s.
        \]
        For all sufficiently large $i$ and every $w = F_i^a/H_{i, w} \in \mathcal{W}_i$, every subring of $\mathbb{Q}_p(t)$ containing $A_s$ and $w$ also contains $g^a/p^i$.
    \end{enumerate}
\end{proposition}

If an element of a subring $T \subseteq D$ has arbitrarily long decompositions in $T$, Lemma~\ref{lem:subring-divisors} supplies divisors $w_i \in \mathcal{W}_i$ in $T$ at arbitrarily large stages.
In Proposition~\ref{prop:subring-rational-algebra}, their poles in part~(2) show that $B$ is finite over the algebra obtained by adjoining the rational constants to $T$.
Part~(3a) separates the retained points, and part~(3b) gives a simple zero of $w_i - w_i(\lambda)$ at each retained point $\lambda$ for all sufficiently large stages in the sequence.
Together these conditions show that this algebra equals $B$.
In Lemma~\ref{lem:subring-local-conductor}, part~(4) forces a power of the two-point evaluation kernel into the local conductor, and the derivative bound in part~(3c) then removes the repeated factors.

With $B_i$ fixed, we begin the choice of $F_i$ by prescribing its values at the two points over one assigned prime $p$.
If $U_i$ must have valuation $h$ at $\gamma \in \{\alpha_p, \beta_p\}$, a target value for $F_i(\gamma)$ is $p^h B_i(\gamma)$.
Indeed, $B_i(\gamma) \ne 0$ by transcendence, so dividing this target by $B_i(\gamma)$ gives $p^h$.
At a $q_n$ stage, if $h = 0$ and $a_n(\gamma)b_{n + 1}(\gamma)$ is a unit, we use $B_i(\gamma)a_n(\gamma)b_{n + 1}(\gamma)$ as the target so that the quotient $a_{n + 1}(\gamma)$ will be close to $1$.
We first construct a polynomial over $\mathbb{Q}_p$ attaining both targets exactly, with the constant term still free to vary.

\begin{lemma}[Local interpolation]
    \label{lem:local-interpolation}
    Let $\alpha, \beta \in \mathbb{Q}_p$ be distinct and nonzero, let $u, v \in \mathbb{Q}_p$, and fix an integer $d \ge 3$.
    Write $L(t) = (t - \alpha)(t - \beta)$.
    There is a monic polynomial $J \in \mathbb{Q}_p[t]$ of degree $d$ such that for every $C \in \mathbb{Q}_p$ the polynomial
    \[
        P_C(t) = J(t) + C \frac{L(t)}{L(0)}
    \]
    is monic of degree $d$ and satisfies $P_C(0) = C$, $P_C(\alpha) = u$, and $P_C(\beta) = v$.
    At either point $\gamma \in \{\alpha, \beta\}$, the derivative $P_C'(\gamma)$ is an affine function of $C$ with nonzero slope.
\end{lemma}

\begin{proof}
    Since $\alpha \ne \beta$, there is a unique polynomial $A$ of degree at most one satisfying $A(\alpha) = u - \alpha^d$ and $A(\beta) = v - \beta^d$.
    Thus $t^d + A(t)$ has the prescribed values at $\alpha$ and $\beta$.
    Since $L$ vanishes at both points and $L(0) = \alpha\beta \ne 0$, we can give this polynomial constant term zero while preserving those values by defining
    \[
        J(t) = t^d + A(t) - A(0)\frac{L(t)}{L(0)}.
    \]
    Adding $C L(t)/L(0)$ now gives $P_C$ the desired constant term and preserves the two evaluations.
    Both corrections have degree at most two, so the leading term remains $t^d$.
    Differentiating gives
    \[
        P_C'(\gamma) = J'(\gamma) + C \frac{L'(\gamma)}{L(0)}.
    \]
    The slope is nonzero because $\alpha$ and $\beta$ are simple roots of $L$.
\end{proof}

For $b_0 = F_1$, the targets at the two points over $2$ are $8$ and $1$.
Lemma~\ref{lem:local-interpolation} gives a local polynomial with $P_C(\alpha_2) = 8$ and $P_C(\beta_2) = 1$ for any chosen constant term $C$.
For a nonconstant $F_1 \in \mathbb{Q}[t]$, the value $F_1(\alpha_2)$ is transcendental over $\mathbb{Q}$, so the exact target $8$ is unattainable.
The congruence $F_1(\alpha_2) \equiv 8 \pmod{16}$ gives $v_2(F_1(\alpha_2)) = 3$, and $F_1(\beta_2) \equiv 1 \pmod{2}$ gives $v_2(F_1(\beta_2)) = 0$.

To pass from the local polynomials to a single rational polynomial, we approximate their coefficients at finitely many primes, as in the adelic approximation theorem of Chabert and Peruginelli \cite[Theorem~3.3 and Remark~3.4]{CPAdelic}.
We also need irreducibility and finitely many nonvanishing conditions on the coefficients.

\begin{lemma}[Choosing rational coefficients]
    \label{lem:rational-coefficients}
    Let $S$ be a finite set of primes, let $\ell$ be a prime outside $S$, and write $R = \mathbb{Z}[1/p \mid p \in S]$.
    Fix an integer $d \ge 2$ and a constant $C \in R$ with $C \equiv (-1)^d \pmod{\ell}$.
    For each $p \in S$ prescribe a nonempty open set $U_p \subseteq \mathbb{Q}_p^{d - 1}$, and let $G_1, \ldots, G_m$ be nonzero polynomials in $d - 1$ variables over $\overline{\mathbb{Q}}$.
    There is a monic polynomial
    \[
        P(t) = t^d + \sum_{j = 1}^{d - 1} c_j t^j + C \in R[t]
    \]
    whose coefficient vector $c = (c_1, \ldots, c_{d - 1})$ belongs to every $U_p$ and satisfies $G_h(c) \ne 0$ for every $h$.
    It may be chosen with $P(t) \equiv (t - 1)^d \pmod{\ell}$ and $P(1) \equiv \ell \pmod{\ell^2}$, so $P(t + 1)$ is Eisenstein at $\ell$ and $P$ is irreducible over $\mathbb{Q}$.
\end{lemma}

\begin{proof}
    Since $C \equiv (-1)^d \pmod{\ell}$, we can prescribe $c_1, \ldots, c_{d - 1}$ modulo $\ell$ to give $P(t) \equiv (t - 1)^d \pmod{\ell}$.
    This gives $P(1) \equiv 0 \pmod{\ell}$.
    Lift the residues to $\ell^2$ and adjust $c_1$ by a multiple of $\ell$ to obtain $P(1) \equiv \ell \pmod{\ell^2}$.
    Thus $P(t + 1) \equiv t^d \pmod{\ell}$ and $v_\ell(P(1)) = 1$, as required for Eisenstein's criterion.

    For each $p \in S$, choose a vector $u_p \in U_p$ and an integer $k_p \ge 1$ with $u_p + p^{k_p}\mathbb{Z}_p^{d - 1} \subseteq U_p$.
    Choose a positive integer $b$ with prime factors in $S$ such that $b u_p \in \mathbb{Z}_p^{d - 1}$ for every $p \in S$.
    Writing $c_j = x_j/b$, we require $x_j \equiv b u_{p, j} \pmod{p^{k_p + v_p(b)}}$ at each $p$.
    Since $\ell \nmid b$, the chosen residues for $c_j$ also prescribe $x_j$ modulo $\ell^2$.
    The Chinese remainder theorem leaves each $x_j$ free to range independently over an infinite arithmetic progression.

    Each $c_j = x_j/b$ therefore ranges independently over an infinite set.
    A polynomial vanishing on the product of these sets is zero, by the one-variable root bound applied successively in each variable.
    Applying this to the nonzero product $G_1 \cdots G_m$ gives the required choice of $c$.
\end{proof}

\begin{proof}[Proof of Proposition~\ref{prop:simultaneous-realization}]
    At stage $i$, we first choose the denominator, the degree, and the primary and guard points.
    We then construct local polynomials at the assigned primes and use their coefficients to choose a single $F_i \in \mathbb{Q}[t]$ satisfying both the valuation conditions and the conditions at retained points.

    \proofheading{The denominator and degree}
    Suppose the stages before $i$ are complete.
    For each exponent tuple in the definition of $\mathcal{W}_i$, write
    \[
        \prod_{h = 1}^{i - 1} U_h^{k_h} = \prod_{h = 1}^{i - 1} F_h^{m_h}.
    \]
    The integers $m_h$ are already determined, since each $U_h = F_h/B_h$ and $B_h$ is a product of powers of the polynomials preceding $F_h$.
    Only finitely many tuples occur, so choose a positive integer $K_i$ larger than every $|m_h|$ that occurs.
    Define $B_i = (F_1 \cdots F_{i - 1})^{K_i}$.
    Substituting $U_i = F_i/B_i$ into the product $w$ gives
    \[
        w = \frac{F_i^a}{H_{i, w}}, \qquad H_{i, w} = \prod_{h = 1}^{i - 1} F_h^{a K_i - m_h}.
    \]
    Since $a \ge 1$, each exponent $a K_i - m_h$ is positive.
    At the first stage the empty product gives $B_1 = H_{1, w} = 1$.

    With these denominators fixed, choose an integer $d_i \ge 5$ greater than every earlier degree $d_h$ and large enough such that $a d_i > \deg H_{i, w}$ for every product on the list.
    We will construct $F_i$ of degree $d_i$.
    Once $F_i$ is irreducible, its degree makes it relatively prime to every earlier $F_h$, so the roots of $H_{i, w}$ are precisely the finite poles of $w$.
    The inequality $a d_i > \deg H_{i, w}$ gives the pole at infinity.

    \proofheading{The primes and their primary and guard points}
    Choose an integer cutoff $M_i > 2^{4 + d_1 + \cdots + d_i}$, with $M_i > M_{i - 1}$ when $i > 1$.
    Let $S_i$ consist of all primes $p \le M_i$, all primes whose evaluation points were assigned at earlier stages, and all primes where the prescribed primary value of $U_i$ is positive.
    Only finitely many primes have been assigned so far, and the prescribed primary value is positive at one prime for $b_n$ and at three primes for $q_n$, so $S_i$ is finite.
    When we choose the coefficients of $F_i$, we will require $F_i \in \mathbb{Z}_p[t]$ and $v_p(F_i(0)) = 0$ for every $p \notin S_i$.
    The increasing cutoffs ensure that every prime eventually belongs to some $S_i$.
    The bound on $M_i$ will give us enough residue classes to avoid the roots of $F_1, \ldots, F_i$ modulo $p$ when we choose guard points at new primes $p > M_i$.

    We now choose the primary point $\alpha_p$ and the guard point $\beta_p$ for each prime $p \in S_i$ whose points have yet to be assigned.
    At this stage we also choose both points for one auxiliary prime outside $S_i$.
    Once chosen, the two points over a prime are fixed for all subsequent stages.

    Fix a newly assigned prime $p \in S_i$ and choose the primary point $\alpha_p \in \mathbb{Z}_p$ transcendental over $\mathbb{Q}$ with $v_p(\alpha_p) = 1$ for odd $p$ and $v_2(\alpha_2) = 2$ for $p = 2$.
    Each prescribed valuation class is uncountable, whereas only countably many elements of $\mathbb{Q}_p$ are algebraic over $\mathbb{Q}$, so such a point exists.
    Since $p$ lies outside every earlier $S_h$, each $F_h$ with $h < i$ has $p$-integral coefficients and $v_p(F_h(0)) = 0$.
    Since $\alpha_p \equiv 0 \pmod{p}$, we have $F_h(\alpha_p) \equiv F_h(0) \pmod{p}$, so $v_p(F_h(\alpha_p)) = 0$.

    To choose the guard point $\beta_p$, first choose its residue modulo $p$, represented by an integer $w_p$.
    We require this residue to be nonzero and to avoid the roots of all earlier $F_h$ modulo $p$.
    At the first stage any $w_p \in \{1, \ldots, p - 1\}$ will do.
    At a later stage we have $p > M_{i - 1}$, since all primes up to $M_{i - 1}$ have already been assigned.
    Peruginelli uses the residue of $\lfloor \log p \rfloor$ to make every fixed nonzero integer polynomial a unit at all sufficiently large primes \cite[proof of Theorem~3.1]{Peruginelli}.
    We allow a range of residues so that we can also avoid the roots of the earlier $F_h$.
    Write $u = \lfloor \log_2 p \rfloor$ and restrict the choice of $w_p$ to the integers from $u$ through $2u$.
    These give $u + 1$ distinct nonzero residues since $2u < p$.
    The cutoff gives $u + 1 > d_1 + \cdots + d_{i - 1}$, while the reductions of the earlier monic polynomials have at most $d_1 + \cdots + d_{i - 1}$ roots altogether.
    We can therefore choose $w_p$ in this interval.
    Choose $\beta_p \in w_p + p\mathbb{Z}_p$ transcendental over $\mathbb{Q}$, which is possible since this residue class is uncountable.
    Our choice of $w_p$ gives $v_p(F_h(\beta_p)) = 0$ for every $h < i$, so $v_p(U_h(\alpha_p)) = v_p(U_h(\beta_p)) = 0$ for all such $h$.

    To verify the assertion about guard valuations in part~(1), fix a nonzero polynomial $V \in \mathbb{Z}[t]$.
    Outside the finite set $S_1$, we have $w_p \ge \lfloor \log_2 p \rfloor$, so $w_p$ tends to infinity with $p$.
    Since $V$ has only finitely many integer roots, $V(w_p) \ne 0$ for all sufficiently large $p$.
    Let $d = \deg V$ and let $C$ be the sum of the absolute values of its coefficients.
    The upper bound $w_p \le 2 \log_2 p$ gives $|V(w_p)| \le C(2 \log_2 p)^d < p$ for all sufficiently large $p$, since $p$ grows faster than any fixed power of $\log_2 p$.
    Thus $V(w_p)$ is a nonzero integer of absolute value less than $p$, so $p \nmid V(w_p)$.
    Since $\beta_p \equiv w_p \pmod{p}$, it follows that $V(\beta_p) \equiv V(w_p) \pmod{p}$ and hence $v_p(V(\beta_p)) = 0$.
    Applying this to the numerator and denominator of a fixed nonzero $f \in \mathbb{Q}(t)$, written as a quotient of polynomials in $\mathbb{Z}[t]$, gives $v_p(f(\beta_p)) = 0$ for all sufficiently large $p$.

    Choose an auxiliary prime $\ell > M_i$ outside $S_i$ at which we will make $F_i(t + 1)$ Eisenstein.
    Choose its primary point $\alpha_\ell$ by the same rule.
    Even at the first stage, take $u = \lfloor \log_2 \ell \rfloor$ and choose the guard representative $w_\ell \in \{u, \ldots, 2u\}$ outside the roots of $F_1, \ldots, F_{i - 1}$ modulo $\ell$.
    The interval contains $u + 1$ distinct residues, and $\ell > M_i$ gives $u + 1 > \sum_{h < i} d_h$, the upper bound on the number of roots to avoid.
    Since $\ell > M_i$, every integer in this interval lies strictly between $1$ and $\ell$, so $w_\ell \not\equiv 1 \pmod{\ell}$.
    Choose a transcendental lift $\beta_\ell$ of this residue.
    We will require $F_i(t) \equiv (t - 1)^{d_i} \pmod{\ell}$, so the chosen residues ensure that $v_\ell(F_i(\alpha_\ell)) = v_\ell(F_i(\beta_\ell)) = 0$.
    Since its points are assigned at stage $i$, the auxiliary prime $\ell$ enters $S_{i + 1}$ as a previously assigned prime.

    \proofheading{The constant term and irreducibility}
    We take the coefficients of $F_i$ in $\mathbb{Z}[1/p \mid p \in S_i]$, so they are integral at every prime outside $S_i$.
    We will use the same rational constant term $C_i = F_i(0)$ in the local polynomials of Lemma~\ref{lem:local-interpolation} at all primes in $S_i$.
    Take $C_i$ of the form
    \[
        C_i = (-1)^{d_i} \prod_{p \in S_i} p^{-N_{p, i}},
    \]
    where each $N_{p, i}$ is a positive multiple of $\ell - 1$.
    Then $v_p(C_i) = -N_{p, i}$ for $p \in S_i$, and $v_q(C_i) = 0$ for every prime $q \notin S_i$.
    Thus we can make $v_p(C_i)$ as negative as needed at any one prime in $S_i$ by increasing its exponent, while preserving the valuations at all other primes.
    We will determine the sizes of these exponents after prescribing the local targets.

    Since $\ell \notin S_i$, Fermat's little theorem gives $p^{\ell - 1} \equiv 1 \pmod{\ell}$ for each $p \in S_i$.
    Hence $C_i \equiv (-1)^{d_i} \pmod{\ell}$, which is the constant term of $(t - 1)^{d_i}$ modulo $\ell$.
    This is the compatibility required in Lemma~\ref{lem:rational-coefficients}, so we can impose $F_i(t) \equiv (t - 1)^{d_i} \pmod{\ell}$ and make $F_i(t + 1)$ Eisenstein at $\ell$.
    Each earlier $F_h$ takes $\ell$-adic unit values at both points over $\ell$, so the denominator $B_i$ does as well.
    Since $F_i$ takes unit values at both points, we obtain $v_\ell(U_i(\alpha_\ell)) = v_\ell(U_i(\beta_\ell)) = 0$, as prescribed.

    \proofheading{Local interpolation at the assigned primes}
    Fix $p \in S_i$.
    For each $\gamma \in \{\alpha_p, \beta_p\}$, let $h$ be the prescribed valuation of $U_i(\gamma)$, so $h = 0$ at the guard.
    The ordinary target for $F_i(\gamma)$ is $\eta_{p, \gamma} = p^h B_i(\gamma)$.
    When $U_i = q_n$, we refine this target at each point where $h = 0$ and $v_p(a_n(\gamma)b_{n + 1}(\gamma)) = 0$ by taking
    \[
        \eta_{p, \gamma} = B_i(\gamma)a_n(\gamma)b_{n + 1}(\gamma).
    \]
    All factors in these targets are already known, and the targets are nonzero because evaluation at either transcendental point is injective on $\mathbb{Q}(t)$.

    Write $L_p(t) = (t - \alpha_p)(t - \beta_p)$.
    Apply Lemma~\ref{lem:local-interpolation} with degree $d_i$ and targets $\eta_{p, \alpha_p}$ and $\eta_{p, \beta_p}$.
    It gives a polynomial $J_{p, i} \in \mathbb{Q}_p[t]$, independent of $C_i$, such that
    \[
        F_{p, i}(t) = J_{p, i}(t) + C_i \frac{L_p(t)}{L_p(0)}
    \]
    is monic of degree $d_i$, has constant term $C_i$, and takes the value $\eta_{p, \gamma}$ at each $\gamma$.

    To preserve the prescribed valuations when choosing the rational polynomial $F_i$, we require
    \[
        \frac{F_i(\gamma)}{\eta_{p, \gamma}} \in 1 + p^r\mathbb{Z}_p,
    \]
    where $r = 1$ for an ordinary target and $r = \max\{1, n\}$ for a refined target at a $q_n$ stage.
    This gives $v_p(U_i(\gamma)) = h$, since $v_p(\eta_{p, \gamma}/B_i(\gamma)) = h$ in both cases.
    At a refined target the recursion gives $a_{n + 1}(\gamma) = \eta_{p, \gamma}/F_i(\gamma)$.
    The inverse of an element of $1 + p^r\mathbb{Z}_p$ belongs to the same group, so $v_p(a_{n + 1}(\gamma) - 1) \ge r$.

    We now choose $N_{p, i}$ large enough for the divisibility estimate in part~(4) and the derivative bound in part~(3).
    For each integer $0 \le s \le i$, express $J_{p, i}$ and the denominators $H_{i, w}$ as polynomials in $p^s t$ over $\mathbb{Q}_p$.
    Only finitely many coefficients occur as $s$ and $w$ vary, so choose an integer $e_p \ge 0$ such that multiplying them all by $p^{e_p}$ gives coefficients in $\mathbb{Z}_p$.
    Since $v_p(1/C_i) = N_{p, i}$, taking $N_{p, i} \ge e_p + i$ gives
    \[
        \frac{J_{p, i}}{C_i} \in p^i A_s, \qquad \frac{H_{i, w}}{C_i} \in p^i A_s,
        \qquad A_s = \mathbb{Z}_p[p^s t],
    \]
    for every $0 \le s \le i$ and every $w \in \mathcal{W}_i$.

    For the derivative bound, fix $w = F_i^a/H_{i, w} \in \mathcal{W}_i$ and write $H = H_{i, w}$.
    Since $F_{p, i}(\gamma) = J_{p, i}(\gamma) = \eta_{p, \gamma}$, the quotient rule gives
    \[
        \left(\frac{F_{p, i}^a}{H}\right)'(\gamma)
        = \left(\frac{J_{p, i}^a}{H}\right)'(\gamma)
          + C_i \frac{a \eta_{p, \gamma}^{a - 1} L_p'(\gamma)}{L_p(0) H(\gamma)}.
    \]
    The coefficient of $C_i$ is nonzero because $\eta_{p, \gamma}$ and $H(\gamma)$ are nonzero and $\gamma$ is a simple root of $L_p$.
    As $N_{p, i}$ increases, the valuation of the second term tends to $-\infty$, while the first term is fixed.
    For sufficiently large $N_{p, i}$, the second term has valuation less than $-i$ and strictly less than the valuation of the first term (with $v_p(0) = +\infty$).
    The valuation of the derivative then equals that of the second term.
    There are only two points $\gamma$ and finitely many products $w$, so we can choose a positive multiple $N_{p, i}$ of $\ell - 1$ satisfying all these bounds and $N_{p, i} \ge e_p + i$.
    We make this choice for each $p \in S_i$ independently, since $v_p(C_i) = -N_{p, i}$ regardless of the exponents at the other primes.

    \proofheading{Choosing $F_i$ over $\mathbb{Q}$}
    We now choose one polynomial that approximates $F_{p, i}$ at every $p \in S_i$ and satisfies the Eisenstein congruences at $\ell$.
    Write a candidate as
    \[
        P(t) = t^{d_i} + \sum_{j = 1}^{d_i - 1} c_j t^j + C_i.
    \]
    At each $p \in S_i$, write $c_{p, j}$ for the coefficient of $t^j$ in $F_{p, i}$.
    If the $c_j$ are sufficiently close to the $c_{p, j}$ in $\mathbb{Q}_p$, then $P(\gamma)/\eta_{p, \gamma} \in 1 + p^r\mathbb{Z}_p$ with the prescribed precision at each of the two points $\gamma$.
    The inequalities $v_p((P^a/H_{i, w})'(\gamma)) < -i$ also persist, since the derivative at $\gamma$ varies continuously with the coefficients and $H_{i, w}(\gamma) \ne 0$.
    For the divisibility estimate we also require
    \[
        \frac{P - F_{p, i}}{C_i} \in p^i A_s \qquad (0 \le s \le i).
    \]
    The coefficient of $(p^s t)^j$ in this quotient is $(c_j - c_{p, j})/(C_i p^{s j})$.
    Thus, for each $p$, we can choose an integer $k_p \ge 1$ such that $v_p(c_j - c_{p, j}) \ge k_p$ for every $j$ suffices for all these finitely many conditions.

    These coefficient neighborhoods give the open sets in Lemma~\ref{lem:rational-coefficients}.
    The conditions at retained points specify the polynomials $G_h$ in that lemma.
    Fix an enumeration of $\overline{\mathbb{Q}}$ and, at stage $i$, take its first $i$ points and discard the roots of $F_1, \ldots, F_{i - 1}$.
    At each remaining point $\lambda$ we require $P(\lambda) \ne 0$.
    For each listed product, write $H = H_{i, w}$ and consider $P^a/H$.
    We require it to take distinct values at every pair of distinct remaining points and to have nonzero derivative at each one.
    Since $H$ is nonzero at all these points, clearing denominators expresses each requirement as the nonvanishing of a polynomial in the free coefficients $c_j$.

    We check that each of these polynomials is nonzero by allowing the coefficients to vary over $\overline{\mathbb{Q}}$.
    We need only vary the coefficients of $t$ and $t^2$, both of which are free since $d_i \ge 5$.
    At zero we already have $P(0) = C_i \ne 0$, and at any nonzero $\lambda$ we can vary $P(\lambda)$ by varying $c_1$.
    For a pair of distinct points, name them $\lambda, \mu$ with $\lambda \ne 0$.
    Adding $x t(t - \mu)$ to $P$, with $x \in \overline{\mathbb{Q}}$, fixes $P(\mu)$ and varies $P(\lambda)$.
    Hence $H(\mu)P(\lambda)^a - H(\lambda)P(\mu)^a$ is a nonconstant polynomial in $x$, so equality of the two values does not hold identically in the $c_j$.

    For the derivative at $\lambda$, take $P(\lambda) \ne 0$.
    The quotient rule shows that $(P^a/H)'(\lambda) \ne 0$ exactly when
    \[
        a H(\lambda)P'(\lambda) - H'(\lambda)P(\lambda) \ne 0.
    \]
    If $\lambda \ne 0$, adding $x t(t - \lambda)$ fixes $P(\lambda)$ and changes $P'(\lambda)$ by $x \lambda$.
    At $\lambda = 0$, adding $x t$ fixes $C_i$ and changes $P'(0)$ by $x$.
    Thus the derivative condition also excludes the zeros of a nonzero polynomial in the $c_j$.

    Lemma~\ref{lem:rational-coefficients} now gives a monic irreducible $F_i \in \mathbb{Z}[1/p \mid p \in S_i][t]$ with constant term $C_i$ satisfying the local and retained-point conditions.
    Each retained point appears among the first $i$ points at every sufficiently late stage and is never discarded, so the separation and nonvanishing requirements give parts~(3a) and~(3b).
    For each fixed $p$, all sufficiently late stages include $p$ in $S_i$, where the bound $v_p(w'(\gamma)) < -i$ for every listed $w$ gives part~(3c).

    \proofheading{The divisibility estimate}
    Fix a prime $p$ and an integer $s \ge 0$, and write $A_s = \mathbb{Z}_p[p^s t]$ as in part~(4).
    For every sufficiently large $i$, we have $p \in S_i$ and $s \le i$, so the coefficient estimates at stage $i$ apply.
    We first compare $g = p^{2s} L_p(t)$ with a scalar multiple of $F_i$.
    The interpolation formula gives
    \[
        \frac{F_i}{C_i} - \frac{L_p}{L_p(0)}
        = \frac{F_i - F_{p, i}}{C_i} + \frac{J_{p, i}}{C_i} \in p^i A_s.
    \]
    The first term on the right belongs to $p^i A_s$ by the choice of the rational coefficients of $F_i$, and the second by the choice of $N_{p, i}$.
    Let $c = p^{2s} L_p(0) \in \mathbb{Z}_p$ and $f_i = c F_i/C_i$.
    Multiplying the displayed difference by $c$ gives $f_i - g \in p^i A_s$.
    In particular, $f_i \in A_s$ since $g \in A_s$.

    Now fix $w = F_i^a/H_{i, w} \in \mathcal{W}_i$.
    Raising $f_i \equiv g \pmod{p^i A_s}$ to the $a$th power gives $u = (g^a - f_i^a)/p^i \in A_s$.
    To express the remaining term $f_i^a/p^i$ as a multiple of $w$, use $F_i^a = H_{i, w}w$ to obtain
    \[
        \frac{f_i^a}{p^i} = v w, \qquad
        v = c^a \left(\frac{1}{C_i}\right)^{a - 1} \frac{H_{i, w}}{p^i C_i}.
    \]
    Here $H_{i, w}/C_i \in p^i A_s$ by the denominator estimate, and $v_p(1/C_i) = N_{p, i} > 0$.
    Since $a \ge 1$, all factors of $v$ belong to $A_s$.
    Thus $g^a/p^i = u + v w$ with $u, v \in A_s$, so every subring of $\mathbb{Q}_p(t)$ containing $A_s$ and $w$ contains $g^a/p^i$.

    \proofheading{The valuation profiles and the congruence for $a_n$}
    Fix a stage $i$.
    At each $p \in S_i$, the target congruences give the prescribed primary valuation of $U_i$ and $v_p(U_i(\beta_p)) = 0$.
    At the auxiliary prime $\ell$, we have verified that both valuations are zero.
    Every other prime outside $S_i$ is assigned at a later stage, when its primary and guard points are chosen so that $F_1, \ldots, F_i$ take unit values at both points.
    Consequently $U_i = F_i/B_i$ also takes unit values there, as required because $S_i$ contains every prime where the prescribed primary valuation is positive.
    Since $U_i$ is a product of integer powers of the polynomials inverted in $B$, its valuation at every polynomial coordinate of the profile is zero.
    This gives the full profiles of the $b_n$ and $q_n$.

    For a fixed nonzero $f \in K$, we have already proved that $v_p(f(\beta_p)) = 0$ for all sufficiently large $p$.
    To compute the primary valuations, write $f = t^k V/W$ with $k = \ord_t(f)$ and $V, W \in \mathbb{Z}[t]$ satisfying $V(0)W(0) \ne 0$.
    For every prime $p \nmid 2V(0)W(0)$, we have $v_p(\alpha_p) = 1$, and $V(\alpha_p)$ and $W(\alpha_p)$ reduce modulo $p$ to the nonzero residues of $V(0)$ and $W(0)$.
    Hence $v_p(f(\alpha_p)) = k$ for all such $p$.

    Starting with $a_0 = t$, each step $a_{n + 1} = a_n b_{n + 1}/q_n$ changes the three primary entries beginning at $p_{2n}$ from $(2, 1, 1)$ to $(0, 0, 2)$ and preserves all the other entries.
    Thus the primary profile of $a_n$ begins with $2n$ zeros, followed by a $2$ and then all ones.
    At every guard we have $v_p(a_n(\beta_p)) = 0$, since the initial value and all factors in the recursion have guard valuation zero.
    Each $a_n$ belongs to $B$ because $t \in B$ and the $b_n$ and $q_n$ are units in $B$.
    The primary and guard valuations are all nonnegative, so $a_n \in D$.

    Fix an integer $m \ge 2$.
    Since $m$ has only finitely many prime divisors, for all sufficiently large $n$ each prime $p = p_j$ dividing $m$ belongs to $S_i$ at the stage $U_i = q_n$, with $j < 2n$ and $v_p(m) \le n$.
    At either point $\gamma \in \{\alpha_p, \beta_p\}$ over such a prime, the profiles give $v_p(a_n(\gamma)) = v_p(b_{n + 1}(\gamma)) = v_p(q_n(\gamma)) = 0$.
    Thus both points received the refined target at that stage, giving
    \[
        v_p(a_{n + 1}(\gamma) - 1) \ge n \ge v_p(m).
    \]
    It follows that $(a_{n + 1} - 1)/m$ takes integral values at both points over each $p \mid m$.
    At every $p \nmid m$ the same holds because $a_{n + 1} \in D$ and $v_p(m) = 0$.
    The quotient belongs to $B$ since $1/m \in \mathbb{Q}$, so the membership test gives $(a_{n + 1} - 1)/m \in D$.
    Shifting the index gives $a_n \equiv 1 \pmod{mD}$ for all sufficiently large $n$.
\end{proof}

\section{Factorization in \texorpdfstring{$D$}{D}}

We prove that $D$ is atomic by induction on the $\ord_t$ value of an element.
At orders zero and one, we will show that each element has only finitely many divisors up to associates.
At higher orders, we will show that $a_n$ divides the element for every sufficiently large $n$ and leaves a cofactor of order one less.

\subsection{Units and the congruence modulo three}

Write $\nu_j(f) = v_{p_j}(f(\alpha_{p_j}))$ for the primary value of $f \in K^\times$ at $p_j$.
Suppose first that every polynomial coordinate of $f$ is zero.
By unique factorization in $\mathbb{Q}[t]$, the only irreducible factors that can occur in its numerator or denominator are the $F_i$.
Thus $f$ is a unit of $B$ and is a rational constant times a finite product of integer powers of the $F_i$.
Substitute $F_i = U_i B_i$, starting with the largest index that occurs.
Since $B_i$ is a product of powers of earlier $F_h$, repeated substitution expresses $f$ as a rational constant $c$ times a finite product of integer powers of the $b_n$ and $q_n$.
For each prime $p$, the evaluations $b_n(\beta_p)$ and $q_n(\beta_p)$ are $p$-adic units for every $n$, so $v_p(f(\beta_p)) = v_p(c)$.
If all guard values of $f$ also vanish, then $c = \pm 1$, so
\begin{equation}
    \label{eq:primary-supported-product}
    f = \pm \prod_{n \ge 0} b_n^{s_n} q_n^{u_n},
\end{equation}
where $s_n, u_n \in \mathbb{Z}$ and only finitely many exponents are nonzero.

For such an $f$, let $g_j = \nu_j(f)$ and set $g_{-1} = u_{-1} = 0$.
At $p_{2k + 1}$ only $q_k$ contributes, whereas at $p_{2k}$ the contributions come from $b_k$, $q_k$, and, when $k > 0$, $q_{k - 1}$.
Their prescribed values give
\[
    g_{2k + 1} = u_k, \qquad
    g_{2k} = 3s_k + 2u_{k - 1} + 2u_k
    \qquad (k \ge 0).
\]
A nonzero element of $D$ is a unit exactly when every coordinate of its profile is zero, since the coordinates of its inverse are their negatives.
For a unit, the odd-indexed primary coordinates therefore give $u_k = 0$ for all $k$, and the even-indexed ones then give $s_k = 0$ for all $k$.
This proves $D^\times = \{1, -1\}$.
It also follows from the divisibility test that two nonzero elements of $D$ with the same profile are associates.

For any $f \in K^\times$ whose guard and polynomial values vanish, the sum of three consecutive primary values centered at $2k$ is $3(s_k + u_{k - 1} + u_k)$.
In particular,
\begin{equation}
    \label{eq:primary-triples}
    g_{2k - 1} + g_{2k} + g_{2k + 1} \equiv 0 \pmod{3}
    \qquad (k \ge 0).
\end{equation}
Suppose $g_j \in \{0, 1\}$ for every $j \ge 2N - 1$, where $N \ge 1$.
For $k \ge N$ the sum in \eqref{eq:primary-triples} is either zero or three, so the three entries are equal.
The triple centered at $2k$ and the next one share the entry $g_{2k + 1}$.
Thus all entries from $2N - 1$ onward are equal, and they are zero because the product for $f$ has finite primary support.

\subsection{The $a_n$ are atoms}

\begin{lemma}
    \label{lem:atoms-an}
    For every $n \ge 0$, the element $a_n$ is an atom of $D$.
\end{lemma}

\begin{proof}
    Since $\ord_t(a_n) = 1$, the element $a_n$ is a nonunit.
    Suppose $a_n = yz$ with $y, z \in D$.
    Their orders at $t = 0$ are nonnegative and sum to one, so we may take $\ord_t(y) = 0$.
    Every guard and every polynomial coordinate other than $\ord_t$ vanishes on $a_n$, hence also on its divisor $y$.
    We can therefore apply \eqref{eq:primary-triples} to $g_j = \nu_j(y)$.

    The bounds $0 \le g_j \le \nu_j(a_n)$ give $g_j = 0$ for $j < 2n$, $g_{2n} \le 2$, and $g_j \in \{0, 1\}$ for $j \ge 2n + 1$.
    Taking $N = n + 1$ in the preceding argument gives $g_j = 0$ for every $j \ge 2n + 1$.
    Now \eqref{eq:primary-triples} at $k = n$ gives $g_{2n} \equiv 0 \pmod{3}$, since both neighbors are zero (with $g_{-1} = 0$ when $n = 0$).
    Hence $g_{2n} = 0$ as well, and $y$ is a unit.
\end{proof}

\subsection{Finitely many divisors at orders zero and one}

\begin{lemma}
    \label{lem:low-order-divisors}
    If $x \in D$ is nonzero and $\ord_t(x) \le 1$, then $x$ has only finitely many divisors up to associates.
\end{lemma}

\begin{proof}
    Only finitely many polynomial coordinates of $x$ are nonzero, since its numerator and denominator have finitely many irreducible factors.
    By Proposition~\ref{prop:simultaneous-realization}(1), all but finitely many guard values are zero, and $\nu_j(x) = \ord_t(x)$ for all sufficiently large $j$.
    If $\ord_t(x) = 0$, the entire profile therefore has finite support.
    At each coordinate the value of a divisor lies between zero and the value of $x$.
    These integer bounds allow only finitely many profiles, and divisors with the same profile are associates.

    Suppose $\ord_t(x) = 1$.
    Now $\nu_j(x) = 1$ for all sufficiently large $j$, so the coordinate bounds alone leave infinitely many places where a divisor might have value one.
    Consider first the divisors $z$ with $\ord_t(z) = 0$.
    Group them by their guard and polynomial values.
    Only finitely many groups occur, since at these coordinates $x$ still has finite support and bounds the values of $z$.

    Fix a representative $u$ of one group.
    For each $z$ in that group, $f = z/u \in K^\times$ has value zero at every guard and polynomial coordinate.
    Choose $N \ge 1$ so that $\nu_j(x) = 1$ and $\nu_j(u) = 0$ for every $j \ge 2N - 1$.
    The same $N$ works for the whole group, since we chose it using only $x$ and $u$.
    For every such $j$ we have
    \[
        \nu_j(f) = \nu_j(z) \in \{0, 1\}.
    \]
    The overlapping triples in \eqref{eq:primary-triples} then give $\nu_j(z) = 0$ throughout this range.
    At the finitely many earlier indices we have $0 \le \nu_j(z) \le \nu_j(x)$, so only finitely many profiles occur within the group.
    This proves finiteness up to associates for the order-zero divisors.
    Every remaining divisor has order one, and $z \mapsto x/z$ pairs its associate class with that of an order-zero divisor.
\end{proof}

Every nonunit $x$ covered by Lemma~\ref{lem:low-order-divisors} factors into atoms.
Indeed, in any decomposition $x = y_1 \cdots y_\ell$ into nonunits, the successive partial products $1, y_1, y_1 y_2, \ldots, x$ are pairwise nonassociate divisors of $x$.
They are distinct; if two were associates, the intervening product of nonunits would be a unit.
Thus the number of divisor classes bounds $\ell + 1$.
In a decomposition of greatest length every factor is an atom, since splitting a reducible factor would increase the length.

\subsection{Atomicity by induction on the order at zero}

\begin{proposition}
    \label{prop:atomicity-D}
    The domain $D$ is atomic.
    In every factorization of a nonzero nonunit $x \in D$ into atoms, exactly $\ord_t(x)$ of the factors, counted with multiplicity, have order one at $t = 0$, and all the others have order zero.
\end{proposition}

\begin{proof}
    Let $r = \ord_t(x)$.
    We have proved atomicity for elements of orders zero and one, so suppose $r \ge 2$.
    By Proposition~\ref{prop:simultaneous-realization}(1), there is an index $J$ such that $\nu_j(x) = r$ for every $j \ge J$.
    Take any $n$ with $2n \ge J$.
    For $j < 2n$ we have $\nu_j(a_n) = 0$, and for $j \ge 2n$ we have $\nu_j(a_n) \le 2 \le r = \nu_j(x)$.
    At $t = 0$ we have $\ord_t(a_n) = 1 \le r$, and all other polynomial and guard values of $a_n$ are zero.
    By coordinatewise comparison, $a_n \mid x$ in $D$ for every sufficiently large $n$.

    Since $\ord_t(x/a_n) = r - 1$, induction gives a factorization of $x/a_n$ into atoms.
    Adjoining the atom $a_n$ gives one for $x$.
    This also shows that every element of order at least two is reducible, since both $a_n$ and $x/a_n$ have positive order.
    Each atom thus has order zero or one, and additivity of $\ord_t$ gives the asserted count.
\end{proof}

The chain beginning at $(t^2)$ now gives a single element with atomic factorizations of unbounded length.
For every $k \ge 1$, the disjoint primary supports of $q_0, q_2, \ldots, q_{2k - 2}$ give
\[
    t^2 = \left(\prod_{j = 0}^{k - 1} q_{2j}\right) x_k,
    \qquad x_k = \frac{t^2}{\prod_{j = 0}^{k - 1} q_{2j}} \in D.
\]
Each $q_{2j}$ is a nonunit of order zero, so $\ord_t(x_k) = 2$.
Factoring the $k$ nonunits and their cofactor into atoms gives a factorization of length at least $k + 2$, since $x_k$ contributes two atoms of order one.
Yet $t = a_0$ is itself an atom, so $t \cdot t$ is a factorization of the same element of length two.

\section{Factorization in every subring}

Fix a subring $T \subseteq D$.
Since $D^\times = \{1, -1\}$ and both units belong to $T$, a nonunit of $T$ remains a nonunit of $D$.
If the lengths of the decompositions of a given nonzero nonunit are bounded, a decomposition of greatest length is a factorization into atoms of $T$.
It remains to consider a subring $T$ in which some nonzero $x$ has arbitrarily long decompositions.
We will use this one element to find a positive integer $m$ with $mD \subseteq T$.
The congruences $a_n \equiv 1 \pmod{mD}$ will then allow us to factor every nonzero nonunit of $T$.

\subsection{Divisors in long decompositions}

\begin{lemma}
    \label{lem:subring-divisors}
    Let $T$ be a subring of $D$ and let $x \in T$ be nonzero with arbitrarily long decompositions in $T$.
    There is a positive integer $R$, depending only on $x$, with the following property.
    For every cutoff $I \ge 1$, one can choose $i > I$ and $w_i \in \mathcal{W}_i \cap T$ such that $x/w_i \in T$ and the exponent of $U_i$ in $w_i$ lies between $1$ and $R$.
\end{lemma}

\begin{proof}
    Let $r = \max_{j \ge 0} \nu_j(x)$, which exists because $\nu_j(x)$ is eventually $\ord_t(x)$, and write $R = 2r$.
    Fix a cutoff $I \ge 1$.
    Each $U_h$ has finite primary support, so we can choose an initial segment of the primary indices containing the supports of all $U_h$ with $h \le \max\{I, R\}$.
    Let $\mathcal{V}$ consist of the primary valuations in this segment, $\ord_t$, and all guard and polynomial valuations at which $x$ has positive value.
    This set is finite by Proposition~\ref{prop:simultaneous-realization}(1) for the guards and by polynomial factorization for the polynomial valuations.
    In any decomposition of $x$, each factor has nonnegative value at every valuation in $\mathcal{V}$.
    At most $\sum_{v \in \mathcal{V}} v(x)$ factors can have positive value anywhere in $\mathcal{V}$, since each such factor contributes at least one to the sum.
    Choose a longer decomposition, and let $y$ be a factor whose value is zero throughout $\mathcal{V}$.
    All remaining guard and polynomial values of $y$ are zero as well, since $y$ divides $x$ and the corresponding values of $x$ are zero.
    By Proposition~\ref{prop:simultaneous-realization}(1), $\nu_j(y) = \ord_t(y) = 0$ for all sufficiently large $j$, so the primary support of $y$ is finite.
    This support is nonempty because $y$ is a nonunit of $D$, and it lies beyond the chosen initial segment because $y$ has value zero throughout that segment.

    By \eqref{eq:primary-supported-product}, write $y = \pm \prod_{n \ge 0} b_n^{s_n} q_n^{u_n}$ with only finitely many nonzero exponents.
    The values $g_j = \nu_j(y)$ satisfy $0 \le g_j \le \nu_j(x) \le r$.
    Since $y$ has nonempty primary support, some $g_j$ is positive, so $r \ge 1$.
    Reading the exponents first at the odd-indexed primary coordinates and then at the even-indexed ones gives
    \[
        u_n = g_{2n + 1}, \qquad 3s_n = g_{2n} - 2u_{n - 1} - 2u_n,
        \qquad u_{-1} = 0.
    \]
    Thus $0 \le u_n \le r$ and $-4r \le 3s_n \le r$, so $R$ bounds the absolute value of every exponent.
    The last generator that occurs in construction order has positive exponent.
    Indeed, if it is $q_n$, then $u_n$ is nonzero and nonnegative.
    If it is $b_n$, then $q_n$ and, for $n \ge 1$, $q_{n - 1}$ both come later, so $3s_n = g_{2n} > 0$.

    The primary profile of $y$ has a positive entry beyond the chosen segment, so its last generator occurs at a stage $i > \max\{I, R\}$.
    Since $R < i$, the exponents satisfy the bounds defining $\mathcal{W}_i$.
    Changing the sign of $y$ if necessary gives the required divisor $w_i$ in $T$.
\end{proof}

Whenever $x$ satisfies the hypothesis of Lemma~\ref{lem:subring-divisors}, choose a sequence of its divisors $w_i$ with strictly increasing stages $i$ as in that lemma.
Write $w_i = F_i^{e_i}/H_{i, w_i}$ as in Proposition~\ref{prop:simultaneous-realization}(2), so $1 \le e_i \le R$ throughout the sequence.

\subsection{Recovering \texorpdfstring{$B$}{B} from the subring}

Let $\mathbb{Q}T$ denote the localization of $T$ at the nonzero integers.
We first recover $t$ after this localization.
The poles of the $w_i$ give a finite extension of rings, and their values and derivatives at the retained points determine its fibres.

\begin{proposition}
    \label{prop:subring-rational-algebra}
    If some nonzero element of $T$ has arbitrarily long decompositions in $T$, then $\mathbb{Q}T = B$.
    In particular, $mt \in T$ for some positive integer $m$.
\end{proposition}

\begin{proof}
    Write $T_{\mathbb{Q}} = \mathbb{Q}T \subseteq B$.
    For each positive integer $h$, choose a divisor $w_i$ from the sequence with $i > h$.
    By Proposition~\ref{prop:simultaneous-realization}(2), its reduced denominator is divisible by $F_h$.
    Write $w_i = P/Q$ in lowest terms, and choose $U, V \in \mathbb{Q}[t]$ with $UP + VQ = 1$.
    Then $1/Q = Uw_i + V \in T_{\mathbb{Q}}[t]$, so $(Q/F_h)/Q = 1/F_h \in T_{\mathbb{Q}}[t]$.
    This gives $B = T_{\mathbb{Q}}[t]$.
    Moreover, $w_i$ has a pole at infinity, so $\deg P > \deg Q$.
    The polynomial $P(X) - w_i Q(X) \in T_{\mathbb{Q}}[X]$ vanishes at $t$ and its leading coefficient is a nonzero rational constant.
    Dividing by that coefficient gives a monic equation for $t$ over $T_{\mathbb{Q}}$.
    Hence $B$ is a finite $T_{\mathbb{Q}}$-module.

    Let $k = \overline{\mathbb{Q}}$, $T_k = k \otimes_{\mathbb{Q}} T_{\mathbb{Q}}$, and $B_k = k \otimes_{\mathbb{Q}} B$.
    Then $T_k \subseteq B_k$ and $B_k$ remains finite over $T_k$.
    The ring $B_k$ is the localization of the principal ideal domain $k[t]$ obtained by inverting the $F_i$.
    The point $0$ is retained, so $(t)$ survives, and the maximal ideals of $B_k$ are precisely $(t - \lambda)B_k$ at the retained points $\lambda$.

    Fix a maximal ideal $\mathfrak{m}$ of $T_k$.
    By the lying over property for module-finite extensions, some $(t - \lambda)B_k$ lies above $\mathfrak{m}$.
    Evaluation at $\lambda$ identifies $T_k/\mathfrak{m}$ with $k$, since its image contains the constants.
    If $\mu$ also lies above $\mathfrak{m}$, then $f - f(\lambda) \in \mathfrak{m}$ gives $f(\mu) = f(\lambda)$ for every $f \in T_k$.
    For distinct retained points a sufficiently late $w_i \in T_{\mathbb{Q}}$ takes different values by Proposition~\ref{prop:simultaneous-realization}(3a).
    Hence $\lambda$ is the unique point over $\mathfrak{m}$.

    To determine the multiplicity at $\lambda$, choose $w_i$ with $w_i'(\lambda) \ne 0$, using part~(3b) of the same proposition.
    Then $w_i - w_i(\lambda) \in \mathfrak{m}$ has a simple zero at $\lambda$.
    The ideal $\mathfrak{m}B_k$ is nonzero, and $(t - \lambda)B_k$ is the only maximal ideal containing it.
    Since $B_k$ is a principal ideal domain, $\mathfrak{m}B_k$ is a positive power of this ideal, and the simple zero forces that power to be one.
    Thus the map $T_k/\mathfrak{m} \to B_k/\mathfrak{m}B_k$ is onto, since both sides are $k$ via evaluation.

    For the finite $T_k$-module $M = B_k/T_k$, this surjectivity says $M/\mathfrak{m}M = 0$.
    Nakayama's lemma gives $M_{\mathfrak{m}} = 0$ at every maximal ideal, whence $M = 0$.
    Since $k$ is faithfully flat over $\mathbb{Q}$, it follows that $B/T_{\mathbb{Q}} = 0$ as well.
    In particular $t \in \mathbb{Q}T$, so clearing an integer denominator gives $mt \in T$ for some $m > 0$.
\end{proof}

\subsection{The conductor at a prime}

Fix a positive integer $m$ with $mt \in T$ and a rational prime $p$.
Under the hypothesis that some nonzero element of $T$ has arbitrarily long decompositions in $T$, we will prove $mD \subseteq T$ for every such $m$.
We first determine the conductor after extending scalars to $\mathbb{Z}_p$.
Write
\[
    T_p = \mathbb{Z}_p \otimes_{\mathbb{Z}} T, \qquad
    D_p = \mathbb{Z}_p \otimes_{\mathbb{Z}} D, \qquad
    B_p = \mathbb{Q}_p \otimes_{\mathbb{Q}} B.
\]
These are ordinary tensor products.
Since $\mathbb{Z}_p$ is torsion-free and hence flat over $\mathbb{Z}$, we can identify $T_p \subseteq D_p \subseteq B_p$.
Here we use $\mathbb{Z}_p \otimes_{\mathbb{Z}} B = B_p$, since $B$ is a $\mathbb{Q}$-algebra.
By Proposition~\ref{prop:subring-rational-algebra}, inverting $p$ gives $T_p[1/p] = B_p$.

The ring $B_p$ is the localization of $\mathbb{Q}_p[t]$ obtained by inverting the $F_i$.
These polynomials have nonzero values at $\alpha_p$ and $\beta_p$ by transcendence over $\mathbb{Q}$, so both evaluations extend to $B_p$.
An element of $D_p$ is a finite $\mathbb{Z}_p$-linear combination of elements of $D$ and therefore has integral values at both points.
Over the coefficient field $\mathbb{Q}_p$ the kernels of the individual evaluations are the distinct maximal ideals $(t - \alpha_p)B_p$ and $(t - \beta_p)B_p$.
Their intersection, the kernel of the pair of evaluations, is $L_p B_p$ for
\[
    L_p(t) = (t - \alpha_p)(t - \beta_p).
\]
Matching the two evaluation values of an element of $mD_p$ with those of an element of $T_p$ leaves an error in $L_p B_p$.
To conclude that the original element belongs to $T_p$, we show that this entire ideal lies in $T_p$.

\begin{lemma}
    \label{lem:subring-local-conductor}
    Suppose some nonzero element of $T$ has arbitrarily long decompositions in $T$.
    The conductor $(T_p : B_p) = \{f \in B_p \mid f B_p \subseteq T_p\}$ equals $L_p B_p$.
\end{lemma}

\begin{proof}
    Write $C = (T_p : B_p)$, the largest ideal of $B_p$ contained in $T_p$.
    Since every element of $B_p = T_p[1/p]$ has the form $u/p^k$ with $u \in T_p$ and $k \ge 0$, we have
    \[
        f \in C \quad\Longleftrightarrow\quad f/p^k \in T_p
        \text{ for every } k \ge 0.
    \]
    We will use this criterion first for a power of $L_p$ and then to remove any repeated factor of a generator of $C$.

    \proofheading{A power of the evaluation kernel}
    Let $s = v_p(m)$.
    Since $m/p^s$ is a unit of $\mathbb{Z}_p$ and $mt \in T_p$, we have $p^s t \in T_p$.
    Thus $T_p$ contains $A_s = \mathbb{Z}_p[p^s t]$ and $g = p^{2s}L_p \in A_s$.
    For each $k \ge 0$, choose a stage $i \ge k$ from our sequence sufficiently late for Proposition~\ref{prop:simultaneous-realization}(4) at $p$ and $s$.
    The subring $T_p$ contains both $A_s$ and $w_i$, so it contains $g^{e_i}/p^i$.
    Using the fixed bound $e_i \le R$ from Lemma~\ref{lem:subring-divisors}, we obtain
    \[
        \frac{g^R}{p^k}
        = p^{i - k} g^{R - e_i} \frac{g^{e_i}}{p^i} \in T_p.
    \]
    The criterion gives $g^R B_p \subseteq C$.
    Since $p$ is a unit in $B_p$, this is $L_p^R B_p \subseteq C$.

    \proofheading{The multiplicities at the two points}
    For $f \in C$ every $f/p^k$ takes integral values at $\alpha_p$ and $\beta_p$.
    Hence each value of $f$ belongs to $\bigcap_{k \ge 0} p^k \mathbb{Z}_p = 0$, giving $C \subseteq L_p B_p$.
    The inclusions $L_p^R B_p \subseteq C \subseteq L_p B_p$ in the principal ideal domain $B_p$ give $C = h B_p$, where $h$ is a product of the two linear factors of $L_p$ with each exponent between $1$ and $R$.

    Suppose $(t - \gamma)^2$ divides $h$ for $\gamma \in \{\alpha_p, \beta_p\}$, and write $q = h/(t - \gamma)^2$.
    We show that $h/(t - \gamma)$ also lies in $C$.
    Choose $d \ge 0$ with $p^d q \in T_p$, using $T_p[1/p] = B_p$.
    Fix $k \ge 0$.
    By Proposition~\ref{prop:simultaneous-realization}(3c), we can choose $w_i$ with $v_p(w_i'(\gamma)) \le -k - d$.
    The scalar $u = 1/(p^k w_i'(\gamma))$ then belongs to $p^d \mathbb{Z}_p$, so $uq \in T_p$.

    Modulo $(t - \gamma)^2 B_p$ the first-order expansion gives
    \[
        w_i(t) - w_i(\gamma) \equiv w_i'(\gamma)(t - \gamma).
    \]
    This expansion is valid for $w_i$ because its denominator is nonzero at $\gamma$ and hence invertible modulo $(t - \gamma)^2$.
    Multiplication by $uq$ sends the error into $h B_p = C$, giving
    \[
        uq\bigl(w_i(t) - w_i(\gamma)\bigr)
        \equiv \frac{h}{p^k(t - \gamma)} \pmod{C}.
    \]
    The left side belongs to $T_p$, since $w_i \in T_p$ and $w_i(\gamma) \in \mathbb{Z}_p$.
    As $C \subseteq T_p$, the right side belongs to $T_p$ for every $k$.
    The criterion now gives $h/(t - \gamma) \in C$.
    This contradicts $C = h B_p$, since $t - \gamma$ is a nonunit of $B_p$.
    Thus both factors have exponent one and $C = L_p B_p$.
\end{proof}

\subsection{An integer in the conductor}

\begin{proposition}
    \label{prop:subring-integer-conductor}
    Suppose some nonzero element of $T$ has arbitrarily long decompositions in $T$.
    Then $mD \subseteq T$ for every positive integer $m$ such that $mt \in T$.
\end{proposition}

\begin{proof}
    Fix such an $m$ and a prime $p$.
    Evaluating $1$ and $t$ at $\alpha_p$ and $\beta_p$ gives the two columns of
    \[
        \begin{pmatrix}
            1 & \alpha_p \\
            1 & \beta_p
        \end{pmatrix}.
    \]
    Its determinant $\beta_p - \alpha_p$ is a unit of $\mathbb{Z}_p$, since $\alpha_p \in p\mathbb{Z}_p$ and $\beta_p \in \mathbb{Z}_p^\times$.
    Thus the two columns form a basis of $\mathbb{Z}_p \times \mathbb{Z}_p$.
    Given $f \in mD_p$, the pair $(f(\alpha_p), f(\beta_p))$ lies in $m(\mathbb{Z}_p \times \mathbb{Z}_p)$ and can be matched by a $\mathbb{Z}_p$-linear combination of $m$ and $mt$.
    This combination belongs to $T_p$, and its difference from $f$ lies in the kernel $L_p B_p \subseteq T_p$ by Lemma~\ref{lem:subring-local-conductor}.
    Hence $mD_p \subseteq T_p$ at every prime $p$.

    To recover the integral inclusion, consider the additive group $M = (mD + T)/T$.
    It is torsion because $\mathbb{Q}T = B$ and thus every element of $D$ has a nonzero integer multiple in $T$.
    By flatness, $\mathbb{Z}_p \otimes_{\mathbb{Z}} M$ identifies with $(mD_p + T_p)/T_p = 0$ for every $p$.
    If $M$ contained a nonzero element of order $e$, then for each prime $p \mid e$, flatness applied to its cyclic subgroup would embed $\mathbb{Z}_p/e\mathbb{Z}_p \ne 0$ in $\mathbb{Z}_p \otimes_{\mathbb{Z}} M$.
    Therefore $M = 0$, which gives $mD \subseteq T$.
\end{proof}

\subsection{Factoring inside the subring}

\begin{proof}[Proof of Theorem~\ref{thm:counterexample}]
    Let $T$ be any subring of $D$.
    If each nonzero nonunit has bounded decomposition lengths in $T$, taking a decomposition of greatest length for each gives atomicity.
    Otherwise Proposition~\ref{prop:subring-rational-algebra} gives a positive integer $m$ with $mt \in T$, and Proposition~\ref{prop:subring-integer-conductor} gives $mD \subseteq T$.
    We can take $m \ge 2$ by passing to a multiple.

    Fix a nonzero nonunit $z \in T$ and let $r = \ord_t(z)$.
    For $r = 0$ or $1$, Lemma~\ref{lem:low-order-divisors} gives only finitely many divisors of $z$ up to associates in $D$.
    The successive partial products in a decomposition of $z$ in $T$ are pairwise nonassociate even in $D$, since each factor is a nonunit of $D$.
    This bounds the lengths of decompositions in $T$, so $z$ factors into atoms of $T$.

    Suppose $r \ge 2$.
    As in the proof of Proposition~\ref{prop:atomicity-D}, every sufficiently late $a_n$ divides $z$ in $D$.
    Choose $n$ large enough that also $a_n \equiv 1 \pmod{mD}$, by Proposition~\ref{prop:simultaneous-realization}(1).
    Then $a_n \in 1 + mD \subseteq T$.
    Write $y = z/a_n \in D$.
    To see that the division takes place in $T$, observe that
    \[
        y - z = y(1 - a_n) \in mD \subseteq T.
    \]
    Thus $y \in T$ as well.
    The element $a_n$ is an atom of $D$ by Lemma~\ref{lem:atoms-an} and remains an atom of $T$ because the unit groups agree.
    Since $\ord_t(y) = r - 1$, induction gives a factorization of $y$, and hence of $z$, into atoms of $T$.

    We have proved that every subring of $D$ is atomic.
    The inclusions $\mathbb{Z}[t] \subseteq D \subseteq \mathbb{Q}(t)$ give $\operatorname{Frac}(D) = \mathbb{Q}(t)$, and the strictly ascending chain beginning at $(t^2)$ proves that $D$ fails ACCP.
\end{proof}

\section{Subrings that fail ACCP}
\label{sec:subring-classification}

Suppose a nonzero element of a subring $T \subseteq D$ has arbitrarily long factorizations into atoms of $T$.
These give arbitrarily long decompositions in $T$, so Propositions~\ref{prop:subring-rational-algebra} and~\ref{prop:subring-integer-conductor} give $mD \subseteq T$ for some positive integer $m$.
We first determine the quotients $D/mD$ and then use them to construct a chain inside each subring of finite additive index.

\subsection{Subrings of finite additive index}

For a positive integer $n$, write $T(n) = \mathbb{Z} + nD$.
The two evaluation points over each prime distinguish the elements of $T(n)$ inside $D$.

\begin{proposition}
    \label{prop:subrings-finite-index}
    For each $n \ge 1$, the ring $T(n)$ is the unique subring of $D$ of additive index $n$.
    The additive quotient $D/T(n)$ is cyclic, generated by the class of $t$, and $(T(n) : D) = nD$.
\end{proposition}

\begin{proof}
    Write $R_n = \mathbb{Z}/n\mathbb{Z}$ for each positive integer $n$.
    We first compute $D/p^hD$ for a prime $p$ and an integer $h \ge 1$.
    For $z \in \mathbb{Z}_p$, let $\bar{z}$ denote its residue in $R_{p^h}$.
    By the definition of $D$, both $f(\alpha_p)$ and $f(\beta_p)$ lie in $\mathbb{Z}_p$ for every $f \in D$, so evaluation and reduction give a ring homomorphism
    \[
        \Phi \colon D \longrightarrow R_{p^h} \times R_{p^h},
        \qquad f \longmapsto \bigl(\overline{f(\alpha_p)}, \overline{f(\beta_p)}\bigr).
    \]

    Every element of $p^hD$ maps to $(0, 0)$.
    Conversely, if $f \in \ker \Phi$, then both values of $f$ lie in $p^h\mathbb{Z}_p$, so $f/p^h$ is integral at both points over $p$.
    For each prime $q \ne p$, division by $p^h$ preserves integrality at $\alpha_q$ and $\beta_q$ because $p$ is a unit in $\mathbb{Z}_q$.
    We also have $f/p^h \in B$ since $1/p^h \in \mathbb{Q} \subseteq B$.
    Thus $f/p^h \in D$, proving $\ker \Phi = p^hD$.

    To obtain a prescribed pair $(u, v)$, it suffices to find $a, b \in \mathbb{Z}$ such that $a + bt$ maps to $(u, v)$.
    Its two residues give the equations
    \[
        \bar{a} + \bar{b}\bar{\alpha}_p = u, \qquad
        \bar{a} + \bar{b}\bar{\beta}_p = v.
    \]
    Since $\alpha_p \in p\mathbb{Z}_p$ and $\beta_p \in \mathbb{Z}_p^\times$, the difference $\bar{\beta}_p - \bar{\alpha}_p$ is a unit in $R_{p^h}$.
    Subtracting the first equation from the second gives $\bar{b} = (v - u)(\bar{\beta}_p - \bar{\alpha}_p)^{-1}$, and then $\bar{a} = u - \bar{b}\bar{\alpha}_p$.
    Taking integer representatives gives the required polynomial $a + bt \in \mathbb{Z}[t] \subseteq D$.
    Hence $\Phi$ is surjective and induces an isomorphism $D/p^hD \cong R_{p^h} \times R_{p^h}$.

    Now suppose $n > 1$.
    Applying the Chinese remainder theorem to the pairwise coprime powers $p^{v_p(n)}$ for $p \mid n$ gives $D/nD \cong R_n \times R_n$, with the primary evaluations over these primes in the first coordinate and the guard evaluations in the second.
    The integers map onto the diagonal, so $T(n)$ is its inverse image.
    In particular, $f \in D$ belongs to $T(n)$ exactly when
    \[
        f(\alpha_p) \equiv f(\beta_p) \pmod{p^{v_p(n)}} \qquad\text{for every } p \mid n.
    \]
    Subtracting the first coordinate from the second identifies the additive quotient $D/T(n)$ with $R_n$.
    The image of $t$ is a unit modulo each $p \mid n$, so its class generates this quotient and $[D : T(n)] = n$.

    If $[D : T] = n$, then $n$ annihilates the additive group $D/T$.
    Hence $nD \subseteq T$ and $T(n) \subseteq T$.
    Since both subgroups have index $n$, they coincide.
    Finally, the conductor $(T(n) : D)$ is the largest ideal of $D$ contained in $T(n)$ and contains $nD$.
    The diagonal in $R_n \times R_n$ contains no nonzero ideal of $R_n \times R_n$, since multiplying a diagonal element $(u, u)$ by $(1, 0)$ gives $(u, 0)$.
    Passing to $D/nD$ therefore gives $(T(n) : D) = nD$.
    When $n = 1$ we have $T(1) = D$, and the assertions follow directly.
\end{proof}

\subsection{Chains in a subring of finite index}

In $D$, successive division by $q_0, q_2, q_4, \ldots$ gives a strict chain beginning at $(t^2)$.
To obtain a chain in $T(n)$, we group sufficiently late terms of the sequence $(q_{2j})_{j \ge 0}$ into products congruent to $1$ modulo $nD$.

\begin{proposition}
    \label{prop:subring-chain-order}
    Let $T$ have finite additive index in $D$ and let $x \in T$ be nonzero.
    The ideal $xT$ starts an infinite strictly ascending chain of principal ideals of $T$ if and only if $\ord_t(x) \ge 2$.
    In that case $x$ has arbitrarily long factorizations into atoms of $T$.
    At orders zero and one, $x$ has only finitely many factorizations into atoms of $T$ up to permutation and replacement of factors by associates.
\end{proposition}

\begin{proof}
    If $\ord_t(x) \le 1$, Lemma~\ref{lem:low-order-divisors} gives finitely many divisors up to associates in $D$.
    Since $T^\times = D^\times = \{1, -1\}$, it gives the same finiteness for divisors in $T$.
    Every principal ideal containing $xT$ is generated by one of these divisors.
    Also, the partial products in a decomposition into nonunits are pairwise nonassociate, so the lengths are bounded.
    With finitely many possible atom factors up to associates, this gives the asserted finiteness of factorizations.

    Suppose $r = \ord_t(x) \ge 2$ and write $T = T(n)$ by Proposition~\ref{prop:subrings-finite-index}.
    Discard enough initial terms of $q_0, q_2, q_4, \ldots$ that $\nu_j(x) = r$ on every remaining primary support and none of those supports lies over a prime dividing $n$.
    This is possible by Proposition~\ref{prop:simultaneous-realization}(1).
    The supports are disjoint and the primary values of each term are at most two.
    Thus any product of distinct remaining terms divides $x$ in $D$.

    If $n > 1$, each remaining term has valuation zero at both evaluations over each $p \mid n$, so its image in $D/nD$ is a unit.
    Among the successive partial products of these terms, some residue occurs infinitely often in the finite group $(D/nD)^\times$.
    Between two successive occurrences of that residue, the product of the intervening terms is congruent to $1$ modulo $nD$.
    In this way we obtain disjoint nonempty blocks with products $C_1, C_2, \ldots$ in $1 + nD \subseteq T$.
    For $n = 1$, take the remaining terms themselves as the $C_i$.
    Each $C_i$ has positive value at a primary and hence is a nonunit of $T$.

    Let $x_0 = x$ and $x_i = x/(C_1 \cdots C_i)$ for $i \ge 1$.
    We know $x_i \in D$, and $x = x_i C_1 \cdots C_i$ gives $x_i \equiv x \pmod{nD}$, so $x_i \in T$.
    From $x_{i - 1} = C_i x_i$ we obtain $x_{i - 1}T \subsetneq x_iT$ at every step.
    Each $x_i$ still has order $r$ at zero, so it too is a nonunit.
    Since $T$ is atomic, factoring $C_1, \ldots, C_i$ and $x_i$ into atoms gives a factorization of $x$ of length at least $i + 1$.
\end{proof}

\begin{proof}[Proof of Theorem~\ref{thm:subring-classification}]
    By Proposition~\ref{prop:subrings-finite-index}, the finite-index subrings are precisely the $T(n)$, with the asserted index and conductor.
    Each $T(n)$ contains $nt^2$, which has order two at zero.
    Proposition~\ref{prop:subring-chain-order} gives both an infinite strictly ascending chain of principal ideals and arbitrarily long factorizations of $nt^2$ into atoms of $T(n)$.

    Conversely, suppose $T$ has infinite additive index.
    Since $D/mD$ is finite for every positive integer $m$, no $mD$ can be contained in $T$.
    By Propositions~\ref{prop:subring-rational-algebra} and~\ref{prop:subring-integer-conductor}, each nonzero nonunit of $T$ has bounded decomposition lengths.
    Thus $T$ has bounded factorization and satisfies ACCP, since a strict infinite chain would give arbitrarily long decompositions of a nonzero generator along the chain.
    Every subring of $T$ also has infinite additive index in $D$, so the same argument gives bounded factorization in each of them.
\end{proof}

In particular, every proper subring of $D$ containing $t$ has infinite additive index, since $t$ generates the additive quotient $D/T(n)$ for every $n > 1$.
Thus it and all its subrings have bounded factorization.

For the polynomial part $A = D \cap \mathbb{Q}[t]$, we can see the infinite index directly.
We have $F_1 = b_0$ with $v_2(F_1(\alpha_2)) = 3$ and value zero at every other coordinate, so $8/F_1 \in D$.
Every nonzero integer multiple of $8/F_1$ still has a pole at the roots of $F_1$, giving an element of infinite order in the additive quotient $D/A$.
Theorem~\ref{thm:subring-classification} therefore gives bounded factorization in every subring of $A$.

For every positive integer $n$, inclusion induces an isomorphism $A/nA \xrightarrow{\sim} D/nD$.
Indeed, the proof of Proposition~\ref{prop:subrings-finite-index} shows that every class in $D/nD$ is represented by an integer linear combination of $1$ and $t$, both of which belong to $A$.
For injectivity, if $f \in A \cap nD$, then $f/n \in D$ is still a polynomial, so $f/n \in A$.
Thus $A$ represents every class modulo $nD$, even though $8n/F_1 \in nD \setminus A$ for every $n$.

\begin{remark}
    Every subring of $A$ is a finite factorization domain.
    To see this, fix a nonzero $f \in A$ and write a divisor of $f$ in $A$ as $cH$, where $c \in \mathbb{Q}^\times$ and $H$ is one of the finitely many monic divisors of $f$ in $\mathbb{Q}[t]$.
    At each guard point we have
    \[
        0 \le v_p\bigl((cH)(\beta_p)\bigr) \le v_p(f(\beta_p)).
    \]
    By Proposition~\ref{prop:simultaneous-realization}(1), the upper bound vanishes for all sufficiently large $p$, so only finitely many tuples of guard values are possible.
    For fixed $H$, subtracting $v_p(H(\beta_p))$ from each entry gives $v_p(c)$ at every prime, which determines $c$ up to sign.
    Thus $f$ has only finitely many divisors in $A$.
    In any subring $T \subseteq A$ containing $f$, each atom occurring in a factorization of $f$ belongs to this finite set.
    Since $T$ has bounded factorization, $f$ has only finitely many factorizations into atoms of $T$ up to order and associates.
\end{remark}

\section{The ideal theory of \texorpdfstring{$D$}{D}}
\label{sec:ideal-theory}

The valuations used to test divisibility also describe the ideals of $D$.
We first identify their centers as the maximal ideals of $D$.
The orders of an ideal at these maximal ideals will then give its factorization into radical ideals and, for invertible ideals, its class in the Picard group.

\subsection{Maximal ideals and the corresponding valuation rings}

For $\gamma \in \{\alpha_p, \beta_p\}$, let $M_{p, \gamma}$ be the kernel of evaluation at $\gamma$ followed by reduction modulo $p$.
Write $P_j = M_{p_j, \alpha_{p_j}}$ and $G_p = M_{p, \beta_p}$ for the maximal ideals corresponding to the primary point $\alpha_{p_j}$ and the guard point $\beta_p$, respectively.
For each monic irreducible $h \in \mathbb{Q}[t]$ other than the $F_i$, let $H_h = D \cap hB$.
In particular $H_t$ consists of the elements vanishing at $t = 0$.

\begin{proposition}
    \label{prop:maximal-ideals-D}
    The maximal ideals of $D$ are the $P_j$, the $G_p$, and the $H_h$.
    The localization at $M_{p, \gamma}$ is the valuation ring of $f \mapsto v_p(f(\gamma))$ on $K$, and $D_{H_h} = \mathbb{Q}[t]_{(h)}$.
    In particular, $D$ is almost Dedekind.
\end{proposition}

\begin{proof}
    By the definition of $A$ and the classification of polynomial overrings in \cite[Theorem~6.2]{CPOverrings}, the closed set of evaluation points for $A$ over $p$ is $E_p = \{\alpha_p, \beta_p\}$.
    Both points are transcendental over $\mathbb{Q}$, so Chang's criterion gives that $A$ is almost Dedekind \cite[Corollary~4.5]{Chang}.
    To pass to $D$, let $M$ be a maximal ideal of $D$ and observe that $A_{M \cap A} \subseteq D_M \subsetneq K$, since $D$ is not a field.
    The contraction is nonzero, since otherwise $D_M$ would contain $\operatorname{Frac}(A) = K$.
    Thus $A_{M \cap A}$ is a discrete valuation ring, and its only strictly larger overring in $K$ is $K$ itself.
    It follows that $D_M = A_{M \cap A}$.

    Over a rational prime $p$, the isomorphism $D/pD \cong \mathbb{F}_p \times \mathbb{F}_p$ from the proof of Proposition~\ref{prop:subrings-finite-index} with $h = 1$ gives exactly the two maximal ideals $M_{p, \alpha_p}$ and $M_{p, \beta_p}$.
    For $M = M_{p, \gamma}$, the evaluation valuation ring contains $D_M$, since each element of $D \setminus M$ has value zero at $\gamma$.
    This valuation ring excludes $1/p$, so it equals the discrete valuation ring $D_M$.

    Inverting the nonzero integers in $D$ gives $B$.
    Indeed, for nonzero $z \in B$, the primary values eventually equal $\ord_t(z) \ge 0$, and the guard values eventually vanish by Proposition~\ref{prop:simultaneous-realization}(1).
    Multiplying $z$ by a suitable positive integer clears the finitely many negative values.
    Since $B$ is obtained from $\mathbb{Q}[t]$ by inverting the $F_i$, its nonzero prime ideals are the $hB$ for the remaining monic irreducibles $h$.
    Their contractions $H_h$ are therefore precisely the nonzero primes of $D$ disjoint from the nonzero integers.
    Each is maximal since $D$ is one-dimensional, and $D_{H_h} = B_{hB} = \mathbb{Q}[t]_{(h)}$.
\end{proof}

Spirito proves that an atomic almost Dedekind domain has zero Jacobson radical whenever it is not Dedekind \cite[Corollary~4.10]{SpiritoBoundness}.
This applies to $D$, since $D$ is atomic and fails ACCP.
For our $D$ the guard points give a direct proof.
Every nonzero $f \in D$ satisfies $v_p(f(\beta_p)) = 0$ for all sufficiently large $p$, so $f$ misses all but finitely many $G_p$.

\subsection{Profiles of invertible ideals}

At each maximal ideal $M$, write $v_M$ for the normalized valuation of $D_M$.
For a nonzero fractional ideal $J$, let $v_M(J)$ be the exponent of $MD_M$ in $JD_M$.
The tuple of these exponents is the \textbf{profile} of $J$.
For a principal fractional ideal $xD$ it is the valuation profile of $x$.
For nonzero integral ideals, this profile is the norm of Coykendall and Hasenauer.
It determines the ideal, and ideal multiplication adds profiles \cite[Definition~2.8 and Theorems~2.9--2.10]{CH}.
For a nonzero fractional ideal $J$, choose nonzero $d \in D$ with $dJ \subseteq D$.
Then $v_M(J) = v_M(dJ) - v_M(d)$ for every maximal ideal $M$.
Write $e_M$ for the tuple with entry one at $M$ and zero elsewhere.

\begin{lemma}
    \label{lem:invertible-profiles}
    An integer tuple $(c_M)_M$ is the profile of an invertible fractional ideal if and only if $c_{P_j} = c_{H_t}$ for all sufficiently large $j$, and only finitely many of its entries at the $G_p$ and at the $H_h$ with $h \ne t$ are nonzero.
\end{lemma}

\begin{proof}
    An invertible fractional ideal has a finite generating set, and its value at $M$ is the minimum of the values of its nonzero generators there.
    Each generator has the stated restrictions by Proposition~\ref{prop:simultaneous-realization}(1) and factorization in $\mathbb{Q}[t]$.
    Taking the minimum over finitely many generators preserves them.

    Conversely, since $D$ is almost Dedekind, it is Pr\"ufer, so each nonzero finitely generated fractional ideal is invertible.
    For each prime $p$, choose an integer $r_p \equiv \beta_p \pmod{p}$.
    The ideals $(p, t)$ and $(p, t - r_p)$ have profiles $e_{M_{p, \alpha_p}}$ and $e_{M_{p, \beta_p}}$, respectively.
    Indeed, $p$ has value one at both points over $p$ and zero at every other coordinate, while $t$ reduces to $0$ at the primary point and to $r_p \not\equiv 0 \pmod{p}$ at the guard.
    For $h \ne t$ as in the definition of $H_h$, the fractional ideal $hD$ has value one at $H_h$ and zero at every other polynomial coordinate.
    Since $\ord_t(h) = 0$, only finitely many of its primary and guard coordinates are nonzero.
    Multiplying by suitable integer powers of the invertible ideals just constructed cancels these entries, realizing $e_{H_h}$.

    Given a tuple satisfying the stated restrictions, subtract $c_{H_t}$ times the profile of $tD$.
    This clears the coordinate at $H_t$ and all but finitely many primary coordinates, since $\nu_j(t) = 1$ for all sufficiently large $j$.
    What remains is a finite integer combination of the tuples $e_M$ with $M \ne H_t$.
    The corresponding product of invertible fractional ideals, multiplied by $t^{c_{H_t}}D$, realizes the original tuple.
\end{proof}

\subsection{$D$ is an SP-domain}

\begin{proposition}
    \label{prop:sp-domain}
    Every nonzero proper ideal of $D$ is a finite product of radical ideals.
    Thus $D$ is an SP-domain.
\end{proposition}

\begin{proof}
    Let $I$ be a nonzero proper ideal of $D$ and write $m_M = v_M(I)$ at each maximal ideal.
    Since $D_M$ is a discrete valuation ring, $m_M$ is the least value of $v_M(x)$ among the nonzero elements $x \in I$.
    A fixed nonzero $x_0 \in I$ gives $0 \le m_M \le v_M(x_0)$ for all $M$.
    The values of $x_0$ are bounded by Lemma~\ref{lem:invertible-profiles} applied to $x_0D$, so $N = \max_M m_M$ is finite.
    Since $I$ is proper, some $m_M$ is positive.
    For each $1 \le k \le N$ we seek a radical ideal with value one exactly where $m_M \ge k$.

    First fix a nonzero $x \in I$.
    The tuple with entry one where $v_M(x) \ge k$ and zero elsewhere satisfies Lemma~\ref{lem:invertible-profiles}, since the primary values of $x$ eventually equal $v_{H_t}(x)$.
    Let $J(x, k)$ be the invertible fractional ideal with this profile.
    Its values are nonnegative, so it is an integral ideal.
    To pass from a single element to all of $I$, define
    \[
        J_k = \sum_{0 \ne x \in I} J(x, k).
    \]
    The sum allows for ideals $I$ that need not be finitely generated.
    Localization commutes with it.
    At $M$, the sum has value one if every nonzero $x \in I$ satisfies $v_M(x) \ge k$, and zero if any summand has value zero.
    Thus $v_M(J_k) = 1$ exactly when $m_M \ge k$.
    Each $J_k$ is radical, since at every maximal ideal it localizes to the whole ring or its maximal ideal.
    At $M$ the product $J_1 \cdots J_N$ therefore has value $m_M$, the number of integers $k$ between one and $m_M$.
    Equality at all maximal localizations gives $I = J_1 \cdots J_N$.
\end{proof}

Spirito identifies the group of invertible fractional ideals of an SP-domain with the group of compactly supported continuous integer-valued functions on its maximal spectrum in the inverse topology, via $J \mapsto (M \mapsto v_M(J))$ \cite[Corollary~4.2]{SpiritoRadical}.
For $D$, this describes the two conditions in Lemma~\ref{lem:invertible-profiles}.
A basic open set consists of the maximal ideals containing a given finitely generated ideal.
In particular, $V(t) = \{H_t\} \cup \{P_j \mid j \ge 0\}$ is open, and $P_j \to H_t$ because the primary values of each generator eventually equal its order at zero.
Every point other than $H_t$ is isolated by an invertible ideal with profile $e_M$.
Thus $V(t)$ is a compact open convergent sequence, and its complement is closed and discrete.
Continuity gives eventual equality of the primary exponents with the exponent at $H_t$, and compact support allows only finitely many nonzero entries outside $V(t)$.

\subsection{The Picard group}

\begin{proposition}
    \label{prop:picard-D}
    The Picard group of $D$ is
    \[
        \operatorname{Pic}(D) \cong \bigoplus_{n \ge 0} \mathbb{Z}/3\mathbb{Z}.
    \]
    The $n$th summand is generated by the class of the invertible ideal with profile $e_{P_{2n}}$.
\end{proposition}

\begin{proof}
    We compute the group of invertible fractional ideals modulo the principal fractional ideals by working with their profiles.
    Start with the profile of an invertible fractional ideal $J$.
    For each $h \ne t$ with $v_{H_h}(J) \ne 0$, subtract the appropriate multiple of the profile of $hD$ to clear that coordinate.
    Next subtract $v_{H_t}(J)$ times the profile of $tD$.
    These steps clear all polynomial coordinates and leave a finitely supported tuple by Lemma~\ref{lem:invertible-profiles}.
    Subtracting multiples of the profiles of the $pD$ clears the guard entries, since the rational prime $p$ has value one at its own guard and zero at every other guard.
    We are left with a finite tuple supported on the primaries.

    Write $e_j = e_{P_j}$.
    The profile of $q_nD$ is $2e_{2n} + e_{2n + 1} + 2e_{2n + 2}$.
    Its coefficient at $2n + 1$ is one, so we can eliminate each odd entry without changing any other odd entry.
    Thus the classes represented by the $e_{2n}$ generate $\operatorname{Pic}(D)$.
    The profile of $b_nD$ is $3e_{2n}$, so the class represented by $e_{2n}$ has order dividing three.

    To determine all relations, suppose a finite sum $\sum_n c_n e_{2n}$ is the profile of a principal fractional ideal $zD$.
    All polynomial and guard values of $z$ vanish, so \eqref{eq:primary-supported-product} gives $z = \pm\prod_n b_n^{s_n}q_n^{u_n}$ with only finitely many nonzero exponents.
    At the primary over $p_{2n + 1}$ its value is $u_n$, so every $u_n$ is zero.
    The even entries are consequently $c_n = 3s_n$ for every $n$.
    This gives precisely the relations $3e_{2n} = 0$ in the quotient, proving the claimed direct sum.
\end{proof}

For comparison, Chang's class group formula gives $\operatorname{Pic}(A) \cong \bigoplus_p \mathbb{Z}$ for the polynomial part of $D$ \cite[Corollary~2.6]{Chang}.
The two evaluation valuations over each $p$ both have $v_p(p) = 1$ and residue field $\mathbb{F}_p$, so the summand at $p$ is $\mathbb{Z}^2/\langle (1, 1) \rangle$.
In $D$, allowing poles at the $F_i$ supplies the principal profiles of the $q_n$ and $b_n$ used in Proposition~\ref{prop:picard-D}.

\section{Variations on the construction}

We outline how to adapt the construction to each positive finite absolute transcendence degree $r$ and each Krull dimension $d$ with $1 \le d \le r + 1$.
We also describe variants that fail weak ACCP or strong atomicity.
The companion Lean library includes formal proofs of these variations \cite{LeanHA}.

\subsection{Krull dimension and absolute transcendence degree}
Fix integers $r \ge 1$ and $1 \le d \le r + 1$.
There is a normal hereditarily atomic domain of dimension $d$ that fails ACCP and has fraction field $\mathbb{Q}(X_1, \ldots, X_r)$.
Its absolute transcendence degree is therefore $r$.
To obtain it, we freshly construct the $F_i$ in all $r$ variables, taking $X_1 = t$ and evaluating at $p$-adic tuples whose coordinates are algebraically independent over $\mathbb{Q}$.
The primary valuations still record the order along $t = 0$, and the $b_n$ and $q_n$ retain their prescribed profiles.

Adjoining an indeterminate $Y$ to $D$ would introduce a nonatomic subring.
Let $z_0 = Yt^2$ and $z_{n + 1} = z_n/q_{2n}$ for $n \ge 0$.
Since $q_0 q_2 \cdots q_{2n}$ divides $t^2$ in $D$ for every $n \ge 0$, each $z_n$ belongs to $YD[Y]$, and the nonunits $q_{2n}$ give a strictly ascending chain $(z_0) \subsetneq (z_1) \subsetneq \cdots$ in $D[Y]$.
Let $C = \mathbb{Z}[q_0, q_2, \ldots] \subseteq D$ and let $P$ be the ideal of $D[Y]$ generated by the $z_n$.
Evaluation at $Y = 0$ fixes $C$ and vanishes on $P$, so $C \cap P = (0)$.
The proof of \cite[Theorem~6.4(6)]{CGH} now shows that $z_0$ has no factorization into atoms in $C + P$, so $C + P$ is a nonatomic subring of $D[Y]$.

The dimension depends on which prime chains survive these choices.
For $d = r$, retain a rational point $c = (1, 0, \ldots, 0)$ in the localization $B$ so that a chain of length $r$ survives there.
The primes over rational primes have height one in this construction.
For $d = r + 1$, also require $f(c) \in \mathbb{Z}_{(q)}$ at one reserved prime $q$.
The kernel of evaluation at $c$ is then strictly contained in the kernel of evaluation followed by reduction modulo $q$, giving one more step.
For $d < r$, preserve a chosen height-$d$ prime and invert an irreducible in every height-$(d + 1)$ prime of the polynomial ring.
A separate guard valuation for each new inverted factor detects its exponent by comparison with the ordinary guard valuation, so the divisibility argument still applies.

\subsection{Two chains witnessing failure of weak ACCP}
\label{sec:two-chains}
A domain $R$ satisfies \emph{weak ACCP} if it is atomic and every nonempty finite set $S$ of nonzero elements of $R$ has a common divisor $e$ such that $x/e$ satisfies ACCP for some $x \in S$ \cite[Proposition~3.5(a) and Definition~3.6]{GL}.
Here an element $z$ satisfies ACCP if every ascending chain of principal ideals beginning at $zR$ stabilizes.

The domain $D$ of Theorem~\ref{thm:counterexample} satisfies weak ACCP.
To see this, fix a nonempty finite set of nonzero elements of $D$.
As long as all the elements have order at least two at zero, a sufficiently late $a_n$ divides them all.
Dividing repeatedly leaves at least one quotient of order zero or one, which has only finitely many divisors up to associates.

We next construct a hereditarily atomic subring of $\mathbb{Q}(t)$ that fails weak ACCP.
Retain a guard point over every rational prime and choose two disjoint infinite families of primes to carry primary points.
For each fixed nonzero $f \in \mathbb{Q}(t)$, the primary valuations of $f$ eventually equal $\ord_t(f)$ in the first family and $\ord_{t - 1}(f)$ in the second.
Interleave the two families of $b_n$ and $q_n$ in a single construction.
The chain from $(t^2)$ now has a second copy starting at $((1 - t)^2)$.
These squares are comaximal, since
\[
    (3 - 2t)t^2 + (1 + 2t)(1 - t)^2 = 1.
\]
Their only common divisors are units, so after division by any common divisor, neither quotient satisfies ACCP.
Thus the set $\{t^2, (1 - t)^2\}$ witnesses the failure of weak ACCP.

\subsection{A pair without a maximal common divisor}
\label{sec:no-maximal-common-divisor}
A \emph{maximal common divisor} of nonzero elements $x, y$ in a domain $R$ is a common divisor $e$ for which the only common divisors of $x/e$ and $y/e$ are units.
The domain $R$ satisfies \emph{strong atomicity} if it is atomic and every pair of nonzero elements has a maximal common divisor \cite[p.~5 and Theorem~1.3(c)]{AAZ}.
Every subring of the two-chain domain in Section~\ref{sec:two-chains} is strongly atomic \cite{LeanHA}.

We now construct a hereditarily atomic domain $D' \subseteq \mathbb{Q}(t)$ containing a pair with no maximal common divisor.
Choose two sequences of $p$-adic evaluation valuations $\nu^+_{k, n}$ and $\nu^-_{k, n}$ for each $k \in \{0, 1\}$, indexed by $n \ge 0$.
Consider a nonzero $f \in D'$ whose profile is supported on these evaluations and the orders at $0$ and $1$.
For both values of $k$, the sequence $\nu^+_{k, n}(f)$ eventually equals $\ord_t(f)$ and the sequence $\nu^-_{k, n}(f)$ eventually equals $\ord_{t - 1}(f)$.
The order of $f$ at $k$ is the weighted sum
\[
    \ord_{t - k}(f) = \sum_{n \ge 0} 2^{-n - 1}\bigl(\nu^+_{k, n}(f) - \nu^-_{k, n}(f)\bigr),
    \qquad k \in \{0, 1\}.
\]
Choose $x$ and $y$ with profiles supported there and with order one at both $0$ and $1$.
For each $k$, the first two $\nu^+$-values are $(2, 0)$ for $x$ and $(0, 4)$ for $y$, and the first two $\nu^-$-values vanish on both.
For $n \ge 2$, all four valuations take value one on both elements.

If $e$ divides both $x$ and $y$, its profile is bounded coordinatewise by both profiles, so the same formula applies to $e$.
Its first two $\nu^+$-values vanish for each $k$, and each later value is at most one.
Thus $0 \le \ord_{t - k}(e) \le \sum_{n \ge 2} 2^{-n - 1} = 1/4$, so both integer orders vanish.
All four sequences of values of $e$ are therefore eventually zero, so its profile has finite support.
Choose $n \ge 2$ beyond that support.
The construction supplies a nonunit with value one at $\nu^+_{0, n}$ and $\nu^-_{0, n}$ and zero elsewhere.
It divides both $x/e$ and $y/e$, so $e$ is not maximal.

Hereditary atomicity comes from a separate feature of these profiles.
Let $T$ be a subring of $D'$ containing $mD'$ for some positive integer $m$, and let $z \in T$ be a nonzero element of positive order at both $0$ and $1$.
The construction gives a factorization $z = uv$ in $D'$ with $\ord_{t - 1}(u) = \ord_t(v) = 0$ and a factorization $m = ab$ into coprime positive integers such that $u \equiv 1 \pmod{aD'}$ and $v \equiv 1 \pmod{bD'}$.
It follows that $u \equiv z \pmod{bD'}$ and $v \equiv z \pmod{aD'}$.
By the Chinese remainder theorem, each factor is congruent modulo $mD'$ to an integer linear combination of $1$ and $z$, so both factors belong to $T$.
Elements of order zero at either point have bounded decomposition lengths, as do all nonzero nonunits in the remaining subrings of $D'$.
Thus every subring is atomic, even though $x$ and $y$ have no maximal common divisor.
This example also fails weak ACCP, since Gotti and Li show that weak ACCP implies strong atomicity \cite[Proposition~3.10]{GL}.

\section{Open questions}

The difficulty in refuting the conjecture is that, to our knowledge, for all previously known unconditional hereditarily atomic domains, the proof establishing hereditary atomicity also establishes hereditary ACCP, thus ruling them out as candidates.
The construction of Theorem~\ref{thm:counterexample} can be adapted to produce a counterexample in characteristic zero in every positive finite absolute transcendence degree, but cannot reach the positive characteristic or absolute transcendence degree zero settings.

\begin{question}
    Does every positive characteristic hereditarily atomic domain satisfy ACCP?
\end{question}

Every subring of $\mathbb{F}_p(t)$ is Noetherian \cite[Theorem~4]{Gilmer}, so a counterexample would require a different fraction field.
For both $R = D$ and $R = D'$, the proof of hereditary atomicity shows that if a nonzero element of a subring $T$ has unbounded decomposition lengths in $T$, then $mR \subseteq T$ for some positive integer $m$.
In a domain $R$ of characteristic $p$, we have $mR = (0)$ if $p \mid m$ and $mR = R$ otherwise, so a proper subring contains no nonzero ideal of this form.

In Theorem~\ref{thm:subring-classification} every subring of $D$ failing ACCP has finite additive index.
In characteristic $p$, Frobenius would identify a counterexample $R$ with its subring $R^p = \{x^p \mid x \in R\}$, giving a subring of infinite additive index that itself fails ACCP.

To see why the index is infinite, suppose $c = \dim_{\mathbb{F}_p}(R/R^p)$ is finite.
For $A = R^{p^n}$, injectivity of Frobenius gives $d = \dim_{\mathbb{F}_p}(R/A) = nc$.
Multiplication embeds $A/I$ in $\operatorname{End}_{\mathbb{F}_p}(R/A)$ for $I = (A : R)$, so $B = R/I$ has dimension at most $d + d^2$.
Choose $n$ large enough that $p^{n - 1} \ge d + d^2$.
Raising to the power $p^{n - 1}$ then kills every nilpotent of $B$, so $B^{p^{n - 1}}$ is finite and reduced and Frobenius maps it onto itself.
Since $I \subseteq A$, the equality $B^{p^{n - 1}} = B^{p^n}$ lifts to $R^{p^{n - 1}} = R^{p^n}$.
Injectivity gives $R = R^p$, although the atomic nonfield $R$ has an atom, which cannot be a $p$th power.

\begin{question}
    Does there exist a hereditarily atomic subring of $\overline{\mathbb{Q}}$ that fails ACCP?
\end{question}

In absolute transcendence degree zero, a hypothetical counterexample $R$ would have to generate an infinite algebraic extension of $\mathbb{Q}$, since every subring of a number field is Noetherian \cite[Theorem~3]{Gilmer}.
Consider the congruence $a_n \equiv 1 \pmod{mD}$ used in the proof of Theorem~\ref{thm:counterexample}.
For each nonzero $x \in R$, a polynomial relation over $\mathbb{Z}$ with nonzero constant term gives a positive integer $N \in xR$.
If a nonunit $a$ divides $x$, then $NR \subseteq aR$, so $a \equiv 1 \pmod{NR}$ would imply $1 \in aR$.
Thus for this single modulus no nonunit divisor of $x$ can play the role of the divisors $a_n$ of $t^2$ in our construction.

The splitting argument in Section~\ref{sec:no-maximal-common-divisor} uses a factorization $m = ab$ with $\gcd(a, b) = 1$, so one factor of the element need only be congruent to $1$ modulo $aD'$ and the other modulo $bD'$.
It relies, however, on the nontrivial valuations $\ord_t$ and $\ord_{t - 1}$, both trivial on $\mathbb{Q}^\times$.
On an algebraic extension of $\mathbb{Q}$, every valuation trivial on $\mathbb{Q}^\times$ is trivial.

There is an alternative hereditary atomicity certificate in $\overline{\mathbb{Q}}$, but no unconditional construction is known to realize it.
A conditional approach uses an infinite tower of totally real number fields, each quadratic over the preceding field, for which the ring of integers of the union is atomic and fails ACCP \cite[Theorem~A]{ConditionalHA}.
This ring can be certified to be hereditarily atomic if every sufficiently late step of the tower satisfies both a Selmer condition at the primes above $2$ and a condition on relative units at the ramified odd primes \cite[Theorem~B]{ConditionalHA}.
The author has not independently verified the conditional construction.
Whether the tower can be chosen to satisfy both conditions at every sufficiently late step remains open.

\section{Acknowledgments}

Under the author's direction, GPT-6.2 Astra discovered the counterexample in Theorem~\ref{thm:counterexample} and the variations thereof.
This construction builds on a six-month investigation by GPT-5.6 Sol, GPT-6 Astra, Claude Opus 5.5, and Claude Fable 5.2.
The author also used AI tools to draft and revise the manuscript.
The author takes full responsibility for the correctness of this manuscript.

The author would also like to thank B.\ Olberding for his suggestions regarding the connections from this construction both to integer-valued polynomials and to D.\ Spirito's work on almost Dedekind domains.

\begin{thebibliography}{99}
    \bibitem{AAZ}
    D.~D. Anderson, D.~F. Anderson, and M.~Zafrullah, \href{https://doi.org/10.1016/0022-4049(90)90074-R}{Factorization in integral domains},
    \emph{J. Pure Appl. Algebra} \textbf{69} (1990), 1--19.

    \bibitem{CPOverrings}
    J.-L. Chabert and G. Peruginelli, \href{https://doi.org/10.1216/JCA-2016-8-1-1}{Polynomial overrings of $\operatorname{Int}(\mathbb{Z})$},
    \emph{J. Commut. Algebra} \textbf{8} (2016), no.~1, 1--28.

    \bibitem{CPAdelic}
    J.-L. Chabert and G. Peruginelli, \href{https://doi.org/10.1016/j.jpaa.2017.04.020}{Adelic versions of the Weierstrass approximation theorem},
    \emph{J. Pure Appl. Algebra} \textbf{222} (2018), no.~3, 568--584.

    \bibitem{Chang}
    G.~W. Chang, \href{https://doi.org/10.4134/JKMS.j210419}{The ideal class group of polynomial overrings of the ring of integers},
    \emph{J. Korean Math. Soc.} \textbf{59} (2022), no.~3, 571--594.

    \bibitem{ConditionalHA}
    \href{https://github.com/mtoenisk/conditional-hereditarily-atomic-without-accp/blob/99c6f3387494fbf02b0840e4393b64b876a90ac1/main.pdf}{\emph{A conditional construction of a hereditarily atomic domain without ACCP}}.
    Manuscript generated with GPT-6 Astra and Claude Opus 5.5, October 22, 2026, commit \texttt{99c6f33}.\newline
    \url{https://github.com/mtoenisk/conditional-hereditarily-atomic-without-accp}.

    \bibitem{CGH}
    J.~Coykendall, F.~Gotti, and R.~E. Hasenauer, \href{https://doi.org/10.1016/j.jpaa.2022.107249}{Hereditary atomicity in integral domains},
    \emph{J. Pure Appl. Algebra} \textbf{227} (2023), no.~4, article~107249.

    \bibitem{CH}
    J.~Coykendall and R.~E. Hasenauer, \href{https://doi.org/10.1017/S0017089517000179}{Factorization in Pr\"ufer domains},
    \emph{Glasg. Math. J.} \textbf{60} (2018), no.~2, 401--409.

    \bibitem{Frisch}
    S.~Frisch, \href{https://doi.org/10.1007/s00605-013-0508-z}{A construction of integer-valued polynomials with prescribed sets of lengths of factorizations},
    \emph{Monatsh. Math.} \textbf{171} (2013), nos.~3--4, 341--350.

    \bibitem{Gilmer}
    R.~Gilmer, \href{https://doi.org/10.1007/BF02567320}{Integral domains with Noetherian subrings},
    \emph{Comment. Math. Helv.} \textbf{45} (1970), 129--134.

    \bibitem{GL}
    F.~Gotti and B.~Li, \href{https://arxiv.org/abs/2212.06213}{Divisibility and a weak ascending chain condition on principal ideals},
    arXiv:2212.06213 (2022).

    \bibitem{GLmonoids}
    F.~Gotti and B.~Li, \href{https://doi.org/10.1007/s00233-025-10578-3}{Arithmetic properties encoded in undermonoids},
    \emph{Semigroup Forum} \textbf{112} (2026), no.~2, 427--449.

    \bibitem{Grams}
    A.~Grams, \href{https://doi.org/10.1017/S0305004100048532}{Atomic rings and the ascending chain condition for principal ideals},
    \emph{Proc. Cambridge Philos. Soc.} \textbf{75} (1974), 321--329.

    \bibitem{LeanHA}
    \emph{Hereditary Atomicity and ACCP}, Lean~4 formalization, version~1.0.0.\newline
    \url{https://github.com/mtoenisk/hereditary-atomicity-and-accp}.

    \bibitem{Peruginelli}
    G.~Peruginelli, \href{https://doi.org/10.2140/pjm.2023.324.333}{Polynomial Dedekind domains with finite residue fields of prime characteristic},
    \emph{Pacific J. Math.} \textbf{324} (2023), no.~2, 333--351.

    \bibitem{SpiritoRadical}
    D. Spirito, \href{https://doi.org/10.1515/forum-2022-0033}{Almost Dedekind domains without radical factorization},
    \emph{Forum Math.} \textbf{35} (2023), no.~2, 363--382.

    \bibitem{SpiritoBoundness}
    D. Spirito, \href{https://doi.org/10.1142/S0219498826501008}{Boundness in almost Dedekind domains},
    \emph{J. Algebra Appl.} \textbf{25} (2026), no.~10, article~2650100.
\end{thebibliography}

\end{document}
